\documentclass[11pt,oneside]{article}

% Standalone ProveIt research note.  The source intentionally has no external
% inputs so that the TeX/PDF pair is directly reproducible after extraction.
\usepackage[a4paper,margin=26mm,headheight=15pt]{geometry}
\usepackage[T1]{fontenc}
\usepackage[utf8]{inputenc}
\IfFileExists{libertinus.sty}{\usepackage{libertinus}}{\usepackage{lmodern}}
\usepackage{microtype}
\usepackage{amsmath,amssymb,amsthm,mathtools,mathrsfs,bm}
\usepackage{booktabs,longtable,tabularx,array}
\usepackage{graphicx}
\usepackage{xcolor}
\usepackage{enumitem}
\usepackage{fancyhdr}
\usepackage{xurl}
\usepackage{listings}
\usepackage[most]{tcolorbox}
\usepackage{hyperref}
\usepackage[nameinlink,noabbrev,capitalise]{cleveref}

\definecolor{linkblue}{RGB}{24,72,124}
\definecolor{deepgreen}{RGB}{23,102,69}
\definecolor{deepred}{RGB}{140,42,42}
\definecolor{deepgold}{RGB}{122,91,19}
\definecolor{shadegray}{RGB}{246,247,249}
\definecolor{rulegray}{RGB}{195,201,208}
\definecolor{palered}{RGB}{251,236,236}
\definecolor{palegold}{RGB}{251,245,230}
\definecolor{palegreen}{RGB}{233,245,243}
\definecolor{paleblue}{RGB}{236,243,251}

\hypersetup{
  hypertexnames=false,
  colorlinks=true,
  linkcolor=linkblue,
  citecolor=deepgreen,
  urlcolor=linkblue,
  pdftitle={Asymptotic Inversion of the Gamma and Barnes G-Functions},
  pdfauthor={Vladimir Reshetnikov},
  pdfsubject={Lambert-core transseries for inverse Gamma and inverse Barnes G, with a general inversion calculus},
  pdfkeywords={inverse gamma function, Barnes G-function, Lambert W, transseries, asymptotic inversion, Hahn series, Stirling series, Lagrange-Burmann inversion}
}

\pagestyle{fancy}
\fancyhf{}
\fancyhead[L]{\small Inverse Gamma and Barnes \(G\)}
\fancyhead[R]{\small\nouppercase{\leftmark}}
\fancyfoot[C]{\thepage}
\renewcommand{\headrulewidth}{0.4pt}

\setcounter{tocdepth}{2}
\setcounter{secnumdepth}{3}
\setlist[itemize]{leftmargin=1.5em,itemsep=0.25ex,topsep=0.6ex}
\setlist[enumerate]{leftmargin=1.9em,itemsep=0.35ex,topsep=0.6ex}
\allowdisplaybreaks[2]

\theoremstyle{plain}
\newtheorem{theorem}{Theorem}[section]
\newtheorem{proposition}[theorem]{Proposition}
\newtheorem{lemma}[theorem]{Lemma}
\newtheorem{corollary}[theorem]{Corollary}
\theoremstyle{definition}
\newtheorem{definition}[theorem]{Definition}
\newtheorem{algorithm}[theorem]{Algorithm}
\newtheorem{example}[theorem]{Example}
\theoremstyle{remark}
\newtheorem{remark}[theorem]{Remark}

\crefname{algorithm}{algorithm}{algorithms}
\Crefname{algorithm}{Algorithm}{Algorithms}

\newtcolorbox{mainbox}[1][]{%
  enhanced,breakable,colback=paleblue,colframe=linkblue,
  boxrule=0.7pt,arc=1.2mm,left=2mm,right=2mm,top=1.4mm,bottom=1.4mm,
  title={Main result},fonttitle=\bfseries,#1}
\newtcolorbox{warningbox}[1][]{%
  enhanced,breakable,colback=palered,colframe=deepred,
  boxrule=0.6pt,arc=1.2mm,left=2mm,right=2mm,top=1.2mm,bottom=1.2mm,
  title={Scale warning},fonttitle=\bfseries,#1}
\newtcolorbox{methodbox}[1][]{%
  enhanced,breakable,colback=palegreen,colframe=deepgreen,
  boxrule=0.6pt,arc=1.2mm,left=2mm,right=2mm,top=1.2mm,bottom=1.2mm,
  title={Reusable method},fonttitle=\bfseries,#1}
\newtcolorbox{summarybox}[1][]{%
  enhanced,breakable,colback=palegold,colframe=deepgold,
  boxrule=0.6pt,arc=1.2mm,left=2mm,right=2mm,top=1.2mm,bottom=1.2mm,
  title={Summary},fonttitle=\bfseries,#1}

\lstdefinestyle{pseudo}{%
  basicstyle=\ttfamily\small,
  columns=fullflexible,
  frame=single,
  rulecolor=\color{rulegray},
  backgroundcolor=\color{shadegray},
  showstringspaces=false,
  keepspaces=true,
  breaklines=true,
  xleftmargin=0.4em,xrightmargin=0.4em,
  aboveskip=0.8em,belowskip=0.8em
}
\lstset{style=pseudo}

\newcommand{\W}{\operatorname{W}}
\newcommand{\e}{\mathrm{e}}
\newcommand{\dd}{\mathop{}\!\mathrm{d}}
\newcommand{\C}{\mathbb{C}}
\newcommand{\R}{\mathbb{R}}
\newcommand{\Q}{\mathbb{Q}}
\newcommand{\N}{\mathbb{N}}
\newcommand{\Tcal}{\mathcal{T}}
\newcommand{\Rcal}{\mathcal{R}}
\newcommand{\Kcal}{\mathcal{K}}
\newcommand{\Bcal}{\mathcal{B}}
\newcommand{\mideal}{\mathfrak{m}}
\newcommand{\coeff}[2]{\left[#1^{#2}\right]}
\newcommand{\littleprec}{\prec}
\DeclareMathOperator{\Log}{Log}
\DeclareMathOperator{\LogGamma}{LogGamma}
\DeclareMathOperator{\InvGamma}{InvGamma}
\DeclareMathOperator{\ord}{ord}
\DeclareMathOperator{\Res}{Res}

\begin{document}

\title{\textbf{Asymptotic Inversion of the Gamma and Barnes \(G\)-Functions}\\[0.35em]
  \Large Lambert-Core Transseries and a General Reversion Calculus}
\author{Vladimir Reshetnikov\\
  \small ProveIt research note}
\date{3 September 2026}
\maketitle

\begin{abstract}
For a large positive value \(y\), this article derives explicit all-orders
asymptotic transseries for the large real solutions of
\(\Gamma(x)=y\) and \(G(x)=y\), where \(G\) is the Barnes \(G\)-function.
The leading inversions are not ordinary power-series reversions.  After taking
logarithms, their dominant phases are respectively
\(x(\log x-1)\) and \(\tfrac12x^2(\log x-\tfrac32)\), and their exact inverse
cores are expressed by the Lambert \(W\)-function.  The lower-order Stirling
terms then generate a second, exponentially finer grid of correction blocks.

For the upper inverse Gamma branch, put
\[
 R=\log\!\left(\frac{y}{\sqrt{2\pi}}\right),\qquad
 w=\W_0(R/\e),\qquad T=\frac{R}{w},\qquad \Lambda=\log T=w+1.
\]
We prove
\[
 \Gamma_+^{-1}(y)\sim \frac12+T+
 \sum_{n\ge1}\frac{d_{2n-1}(\Lambda)}{T^{2n-1}},
\]
where \(d_{2n-1}\in\Lambda^{-1}\Q[\Lambda^{-1}]\), give a triangular
coefficient recurrence, prove remainder transfer from Stirling's series, and
list the first five nonzero blocks through \(T^{-9}\).  The first two blocks
agree with the finite expansion previously proved by Gerhold, Hubalek, and
Tomovski; the recurrence continues it to arbitrary order.  We also record the
ordinary convergent inverse series for the second positive branch tending to
zero.

For Barnes \(G\), it is cleaner to solve \(G(z+1)=y\).  With
\[
 C_G=\zeta'(-1),\quad R=\log y-C_G,\quad
 v=\W_0(4R/\e^3),\quad U=2\sqrt{R/v},\quad
 \lambda=\frac{v+1}{2}=\log U-1,
\]
we obtain
\[
 z\sim U+\sum_{n\ge0}\frac{b_n(\lambda)}{U^n},
 \qquad G_+^{-1}(y)=1+z.
\]
The coefficient functions are explicit rational Laurent polynomials in
\(\lambda\) with coefficients in \(\Q[\tfrac12\log(2\pi)]\).  A logarithmic
term in the forward Barnes expansion is resonant with the core slope; it
forces the nonzero coefficientwise limit
\[
 \sum_{n\ge1}e_n(\lambda)q^n
 \longrightarrow \sqrt{1+q^2/6}-1
 \qquad(\lambda\to\infty),
\]
which explains the limiting constants \(1/12,-1/288,1/3456,\ldots\) in the
odd inverse-power blocks.

The two examples are instances of a general theory for phases
\(a x^p(\log x-\beta)+\Psi(x)\).  We develop exact Lambert-core inversion,
a Hahn-series implicit-function theorem, a support-semigroup calculus,
triangular and Newton--Hensel algorithms, a Lagrange--B\"urmann reduction to
near-identity reversion, denominator bounds, resonant logarithmic sectors,
complex branch labels, residual error certificates, and transport of
exponentially small sectors.  Formal, Poincar\'e-asymptotic, and convergent
claims are kept separate throughout.
\end{abstract}

\noindent\textbf{Keywords.}
Inverse Gamma function; Barnes \(G\)-function; Lambert \(W\); transseries;
series reversion; Hahn series; Stirling series; Lagrange--B\"urmann inversion;
Newton--Hensel lifting; exponentially improved asymptotics.

{\small\tableofcontents}
\clearpage

\section{Introduction and principal formulas}
\label{sec:introduction}

The inversion of a rapidly growing special function has two logically distinct
layers.  The first layer inverts the dominant growth law.  The second layer
propagates every lower-order asymptotic term through the inverse map.  For
polynomial growth the first layer is algebraic; for an exponential it is a
logarithm; for a phase of the form \(x^p\log x\), the exact core inverse is a
Lambert \(W\)-expression.  Gamma and Barnes \(G\) are the first two canonical
examples of the latter class.

The familiar leading estimate for inverse Gamma follows by solving Stirling's
main term.  What is less visible in that one-line calculation is the hierarchy
of scales.  Expanding Lambert \(W\) creates an outer series in powers of
\(1/\log\log y\), with polynomial dependence on \(\log\log\log y\).  A fixed
shift, such as the \(1/2\) in the optimal Gamma centering, is smaller than every
term of that outer scale.  The Bernoulli corrections are smaller again: they
form a grid in powers of the inverse Lambert core, with coefficient functions
that themselves have inverse-logarithmic expansions.  Treating all these
objects as if they were terms of one ordinary power series obscures both the
algebra and the attainable error estimates.

The Barnes problem has one further feature.  Its forward phase contains
\(-\tfrac1{12}\log z\).  After normalization by the dominant
\(\tfrac12z^2\log z\) phase, this term is small by \(z^{-2}\), but its factor
\(\log z\) is of the same order as the core slope parameter.  It therefore
survives in the coefficientwise large-slope limit and creates an algebraic
subseries inside the inverse transseries.  This is the simplest example of the
\emph{resonant logarithmic correction} developed in \cref{sec:resonance}.

\subsection{Branch conventions}

On the positive real axis, \(\Gamma\) has a unique minimum: indeed,
\(\psi'(x)=\sum_{n\ge0}(n+x)^{-2}>0\), while
\(\psi(x)\to-\infty\) as \(x\to0^+\) and
\(\psi(x)\to+\infty\) as \(x\to+\infty\).  Thus \(\psi\) has one
zero and \(\Gamma'=\Gamma\psi\) changes sign exactly once.  Consequently,
a large positive \(y\) has two positive preimages.  We write
\(\Gamma_+^{-1}(y)\) for the branch tending to \(+\infty\) and
\(\Gamma_0^{-1}(y)\) for the branch tending to \(0^+\).  The Lambert-core
transseries concerns \(\Gamma_+^{-1}\); the second branch has an ordinary
series in \(1/y\), given in \cref{sec:gamma-small-branch}.

For Barnes \(G\), we use \(\Bcal(y)\) for the eventual positive real inverse of
\(z\mapsto G(z+1)\):
\[
 G(\Bcal(y)+1)=y,\qquad \Bcal(y)\longrightarrow+\infty.
\]
Thus the corresponding branch of the inverse of \(G\) itself is
\[
 G_+^{-1}(y)=1+\Bcal(y).
\]
This convention removes a persistent off-by-one shift from the asymptotic
formulas.

\subsection{The answer at a glance}

\begin{mainbox}
Let \(B_n(x)\) denote a Bernoulli polynomial and \(B_n=B_n(0)\).

\textbf{Upper inverse Gamma.}  Define
\begin{equation}
 R_\Gamma=\log y-\frac12\log(2\pi),\quad
 w=\W_0(R_\Gamma/\e),\quad
 T=\frac{R_\Gamma}{w},\quad
 \Lambda=\log T=w+1.
 \label{eq:gamma-core-summary}
\end{equation}
Then, for each fixed \(N\ge0\),
\begin{equation}
 \Gamma_+^{-1}(y)
 =\frac12+T+\sum_{n=1}^{N}
   \frac{d_{2n-1}(\Lambda)}{T^{2n-1}}
 +O\!\left(\frac{1}{T^{2N+1}\Lambda}\right),
 \qquad y\to+\infty,
 \label{eq:gamma-summary}
\end{equation}
where the \(d_{2n-1}\) are generated by \cref{eq:gamma-recurrence}.  In
particular,
\begin{align}
 d_1(\Lambda)&=\frac{1}{24\Lambda},\label{eq:gamma-d1-summary}\\
 d_3(\Lambda)&=-\frac{14\Lambda^2+10\Lambda+5}
 {5760\Lambda^3},\label{eq:gamma-d3-summary}\\
 d_5(\Lambda)&=\frac{2232\Lambda^4+1176\Lambda^3+714\Lambda^2
 +350\Lambda+105}{2903040\Lambda^5}.
 \label{eq:gamma-d5-summary}
\end{align}

\textbf{Upper inverse Barnes \(G\).}  Put
\begin{equation}
 \alpha=\frac12\log(2\pi),\quad C_G=\zeta'(-1),\quad
 R_G=\log y-C_G,
 \label{eq:barnes-constants-summary}
\end{equation}
\begin{equation}
 v=\W_0(4R_G/\e^3),\qquad
 U=2\sqrt{R_G/v},\qquad
 \lambda=\frac{v+1}{2}=\log U-1.
 \label{eq:barnes-core-summary}
\end{equation}
Then, for each fixed \(N\ge0\),
\begin{equation}
 \Bcal(y)=U+\sum_{n=0}^{N}\frac{b_n(\lambda)}{U^n}
 +O(U^{-N-1}),
 \qquad y\to+\infty,
 \label{eq:barnes-summary}
\end{equation}
with coefficients generated by \cref{eq:barnes-recurrence}.  The first three
are
\begin{align}
 b_0(\lambda)&=-\frac{\alpha}{\lambda},
 \label{eq:barnes-b0-summary}\\
 b_1(\lambda)&=
 \frac{6\alpha^2\lambda-6\alpha^2+\lambda^3+\lambda^2}
 {12\lambda^3},
 \label{eq:barnes-b1-summary}\\
 b_2(\lambda)&=
 \frac{\alpha(8\alpha^2\lambda-6\alpha^2+\lambda^2)}
 {12\lambda^5}.
 \label{eq:barnes-b2-summary}
\end{align}
Finally, \(G_+^{-1}(y)=1+\Bcal(y)\).
\end{mainbox}

\subsection{Relation to earlier work}

The forward asymptotic expansions used below are classical.  We take the
Gamma form from the NIST Digital Library of Mathematical Functions
\cite{DLMF511} and the Barnes form from DLMF and Nemes
\cite{DLMF517,NemesBarnes}.  Corless, Gonnet, Hare, Jeffrey, and Knuth provide
the standard analytic and asymptotic framework for Lambert \(W\)
\cite{CorlessEtAl}.  Jeffrey and Watt discuss branches of the complex inverse
Gamma function and record the leading Lambert approximation
\cite{JeffreyWatt}.  Gerhold, Hubalek, and Tomovski prove the inverse-Gamma
expansion through the \(T^{-3}\) block in notation equivalent to ours
\cite{GerholdHubalekTomovski}.  The present derivation recovers those terms,
continues them by a closed triangular recurrence, treats Barnes \(G\) in the
same framework, and isolates the general support and resonance mechanisms.

No assertion of convergence is hidden in the symbol \(\sim\).  Stirling's
Bernoulli series is generally divergent, so the inverse block expansions are
Poincar\'e asymptotic transseries.  A finite truncation is rigorous once the
corresponding forward remainder is controlled; \cref{sec:residuals} explains
that transfer.  Exponentially improved forward expansions lead to deeper
inverse sectors, but they do not make the untruncated Bernoulli series
convergent.

\begin{warningbox}
There are three different operations that are easily conflated:
\begin{enumerate}
\item retaining Lambert \(W\) exactly and expanding only the lower-order
      phase corrections;
\item expanding Lambert \(W\) itself in iterated logarithms;
\item optimally truncating the divergent Bernoulli sector and then adding
      exponentially small terms.
\end{enumerate}
They correspond to three nested asymptotic levels.  An equality or error term
valid at one level need not survive a reordering at another.
\end{warningbox}

\section{Exact inversion of the Lambert core}
\label{sec:core}

\subsection{The model phase}

Let
\begin{equation}
 \Phi_0(x)=a x^p(\log x-\beta),
 \qquad a>0,\quad p>0,
 \label{eq:general-core}
\end{equation}
and suppose first that all quantities are positive real and that the branch at
large \(x\) is intended.  The following elementary lemma is the organizing
identity for the whole article.

\begin{lemma}[Exact Lambert-core inverse]
\label{lem:core-inverse}
Let \(R>0\), and define
\begin{equation}
 Z=\frac{pR}{a}\e^{-p\beta},\qquad
 v=\W_0(Z),\qquad
 X=\e^\beta\e^{v/p}
   =\left(\frac{pR}{av}\right)^{1/p}.
 \label{eq:general-core-variables}
\end{equation}
Then \(\Phi_0(X)=R\).  Moreover,
\begin{equation}
 \log X=\beta+\frac{v}{p},\qquad
 \Phi_0'(X)=aX^{p-1}(v+1),
 \label{eq:core-slope}
\end{equation}
and
\begin{equation}
 \frac{\dd X}{\dd R}=\frac{1}{aX^{p-1}(v+1)}.
 \label{eq:core-derivative}
\end{equation}
\end{lemma}

\begin{proof}
Put \(u=p(\log X-\beta)\).  Equation \(\Phi_0(X)=R\) is equivalent to
\[
 \frac{a}{p}\e^{p\beta}u\e^u=R,
\]
or \(u\e^u=Z\).  Thus \(u=v=\W_0(Z)\), which gives
\cref{eq:general-core-variables}.  Differentiating
\(a x^p(\log x-\beta)\) and substituting
\(p(\log X-\beta)=v\) yields \cref{eq:core-slope}; the inverse-function
formula then gives \cref{eq:core-derivative}.
\end{proof}

The derivative in \cref{eq:core-slope} identifies the universal divisor
\(v+1\) that appears in every correction coefficient.  It also explains why a
small phase residual should be divided by \(aX^{p-1}(v+1)\), not merely by the
leading power \(X^{p-1}\).

\subsection{The outer iterated-logarithmic layer}

For \(|Z|\to\infty\) in a fixed sector avoiding the standard cuts, write
\[
 L_1=\log Z,
 \qquad L_2=\log L_1.
\]
The classical expansion
\begin{equation}
 \W(Z)\sim L_1-L_2+\frac{L_2}{L_1}
 +\frac{L_2(-2+L_2)}{2L_1^2}
 +\frac{L_2(6-9L_2+2L_2^2)}{6L_1^3}+\cdots
 \label{eq:W-outer}
\end{equation}
turns \(X=(pR/(av))^{1/p}\) into a polynomial-logarithmic transseries in
\(1/L_1\).  We deliberately keep \(v=\W(Z)\) exact for most of the article.
Doing so packages infinitely many outer logarithmic terms into one analytic
quantity and exposes the much smaller correction grid generated by the
remainder of the phase.

Define
\begin{equation}
 q=X^{-1}=\e^{-\beta-v/p}.
 \label{eq:q-general}
\end{equation}
As \(v\to+\infty\),
\begin{equation}
 q\littleprec v^{-N}\qquad\text{for every fixed }N>0.
 \label{eq:q-beyond-v}
\end{equation}
Thus powers of \(q\) are beyond all orders of the coefficient scale
\(v^{-1}\).  The natural inverse object is consequently a two-level
transseries: a Hahn series in \(q\), whose coefficients are rational or
formal Laurent series in \(v^{-1}\).

\section{Normalized displacement and the triangular inverse theorem}
\label{sec:triangular}

\subsection{The exact normalized equation}

Consider a perturbed phase
\begin{equation}
 \Phi(x)=a x^p(\log x-\beta)+C+\Psi(x)
 \label{eq:perturbed-phase}
\end{equation}
and the equation \(\Phi(x)=Y\).  Set \(R=Y-C\), construct \(v,X,q\) by
\cref{eq:general-core-variables,eq:q-general}, and write
\begin{equation}
 x=X(1+\epsilon).
 \label{eq:relative-displacement}
\end{equation}
After division by \(aX^p\), the core displacement is
\begin{equation}
 \Kcal_{p,v}(\epsilon)
 =(1+\epsilon)^p\left(\frac{v}{p}+\log(1+\epsilon)\right)-\frac{v}{p}.
 \label{eq:kernel-general}
\end{equation}
The normalized perturbation is
\begin{equation}
 \Rcal(q,v,\epsilon)
 =\frac{\Psi(X(1+\epsilon))}{aX^p}.
 \label{eq:normalized-remainder}
\end{equation}
Hence the inverse problem is exactly
\begin{equation}
 \Kcal_{p,v}(\epsilon)+\Rcal(q,v,\epsilon)=0.
 \label{eq:normalized-equation}
\end{equation}

The Taylor coefficients of the universal kernel have a compact closed form:
\begin{equation}
 \Kcal_{p,v}(\epsilon)=\sum_{n\ge1}\kappa_n(p,v)\epsilon^n,
 \qquad
 \kappa_n(p,v)=\frac{v}{p}\binom pn+
 \frac{\partial}{\partial p}\binom pn.
 \label{eq:kappa-general}
\end{equation}
Indeed, \((1+\epsilon)^p\log(1+\epsilon)=\partial_p(1+\epsilon)^p\).
In particular,
\begin{equation}
 \kappa_1(p,v)=v+1,
 \label{eq:kernel-linear}
\end{equation}
which is the normalized core slope.

\subsection{Hahn-series setting}

Let \(S\subset\R_{\ge0}\) be a well-ordered additive submonoid with
\(0\in S\), and let \(\mathscr C\) be a characteristic-zero coefficient
field containing \(p,v\) and \((v+1)^{-1}\).  Typical choices are
\(\C(v)\), \(\C((v^{-1}))\), or a differential transseries field extending
them.  Write
\[
 \mathscr C[[q^S]]
 =\left\{\sum_{\rho\in S}c_\rho q^\rho:
   \operatorname{supp}(c)\text{ is well ordered}\right\}
\]
for the corresponding Hahn ring and
\[
 \mideal=\left\{F\in\mathscr C[[q^S]]:\ord_qF>0\right\}
\]
for its positive-valuation ideal.

The condition that \(S\) be well ordered is not decorative.  It ensures that
at any prescribed exponent \(\rho\), only finitely many decompositions of
\(\rho\) can contribute to the coefficient of a product or composition under
the finite-type hypotheses used below.  For the two principal examples,
\(S=2\N\) for shifted Gamma and \(S=\N\) for Barnes \(G\), so no transfinite
bookkeeping is needed; the Hahn formulation states the closure principle that
survives fractional powers and mixed grids.

\begin{theorem}[Triangular Lambert-core inversion]
\label{thm:triangular}
Assume that
\begin{equation}
 \Rcal(q,v,E)=
 \sum_{\substack{\rho\in S,\ \rho>0\\m\ge0}}
 r_{\rho,m}(v)q^\rho E^m
 \label{eq:R-hahn}
\end{equation}
is locally finite at every \(q\)-exponent and belongs to
\(\mideal\mathscr C[[E]]\).  If \(v+1\) is invertible in \(\mathscr C\), then
\cref{eq:normalized-equation} has a unique solution
\begin{equation}
 \epsilon(q,v)=\sum_{\substack{\gamma\in S\\\gamma>0}}
 e_\gamma(v)q^\gamma\in\mideal.
 \label{eq:epsilon-hahn}
\end{equation}
For every \(\gamma>0\), let
\(\epsilon_{<\gamma}=\sum_{0<\delta<\gamma}e_\delta q^\delta\).  Then
\begin{equation}
 e_\gamma(v)=-\frac{1}{v+1}\coeff{q}{\gamma}
 \left(
   \sum_{n\ge2}\kappa_n(p,v)\epsilon_{<\gamma}^{\,n}
   +\Rcal(q,v,\epsilon_{<\gamma})
 \right).
 \label{eq:general-triangular-recurrence}
\end{equation}
\end{theorem}

\begin{proof}
Because \(\epsilon\in\mideal\), every nonlinear power \(\epsilon^n\),
\(n\ge2\), has strictly larger valuation than \(\epsilon\).  Likewise, every
monomial \(q^\rho\epsilon^m\) in \(\Rcal\) has \(\rho>0\).  Therefore the
coefficient of \(q^\gamma\) in all terms of
\cref{eq:normalized-equation} except \((v+1)\epsilon\) depends only on
coefficients \(e_\delta\) with \(\delta<\gamma\).  Local finiteness makes that
coefficient a finite sum.  Solving the coefficient equation gives
\cref{eq:general-triangular-recurrence}; transfinite induction over the
well-ordered support constructs a solution.

For uniqueness, suppose two solutions first differ at exponent \(\gamma\).
All nonlinear and perturbed contributions at that exponent depend only on
lower coefficients and therefore agree.  The difference of the two equations
is \((v+1)(e_\gamma-\widetilde e_\gamma)=0\), so the coefficients agree, a
contradiction.
\end{proof}

\begin{corollary}[Support semigroup]
\label{cor:support}
If the \(q\)-support of \(\Rcal(q,v,0)\) is contained in a set \(A\subset
\R_{>0}\), and all \(E\)-dependent coefficients of \(\Rcal\) use the same
base exponents, then the support of \(\epsilon\) is contained in the additive
semigroup generated by \(A\).
\end{corollary}

\begin{proof}
The recurrence combines previously occurring exponents only by addition and
adds base exponents from \(A\).  Induction gives the claim.
\end{proof}

\subsection{A denominator bound on an integral grid}

The ubiquitous powers of the slope divisor have a simple combinatorial
origin.

\begin{proposition}[Universal slope-denominator bound]
\label{prop:denominator-bound}
Suppose \(S=\N\), the coefficients of \(\Rcal\) lie in a polynomial ring
\(A[v]\), and no coefficient contains \((v+1)^{-1}\).  If
\(\epsilon=\sum_{n\ge1}e_nq^n\), then
\begin{equation}
 (v+1)^{2n-1}e_n\in A[v]
 \qquad(n\ge1).
 \label{eq:denominator-bound}
\end{equation}
The exponent \(2n-1\) is an upper bound; cancellations may lower it.
\end{proposition}

\begin{proof}
The assertion is immediate for \(n=1\).  Assume it for every lower index.
A coefficient of \(q^n\) in a nonlinear product
\(e_{j_1}\cdots e_{j_m}\), where \(m\ge2\) and
\(j_1+\cdots+j_m=n\), has denominator exponent at most
\[
 \sum_{r=1}^m(2j_r-1)=2n-m\le2n-2.
\]
After the final division by \(v+1\) in
\cref{eq:general-triangular-recurrence}, the exponent is at most \(2n-1\).
A monomial from \(\Rcal\) carries an additional positive base exponent
\(q^\rho\), so its indices sum to \(n-\rho<n\) and gives a strictly smaller
bound.  This closes the induction.
\end{proof}

\subsection{Formal versus analytic meaning}

\Cref{thm:triangular} is a formal implicit-function theorem.  To obtain an
asymptotic theorem for an analytic function, one also needs:
\begin{enumerate}
\item a forward Poincar\'e expansion for \(\Psi\), with a remainder uniform in
      the sector or interval under consideration;
\item eventual nonvanishing of \(\Phi'(x)\);
\item a residual-to-root estimate, supplied in \cref{sec:residuals}.
\end{enumerate}
The formal recurrence then computes a candidate truncation, while the analytic
hypotheses prove that the exact inverse differs from it by the stated next
scale.  This separation prevents a common error: formal existence in a Hahn
ring does not imply convergence of the resulting infinite series.

\section{Equivalent reversion mechanisms}
\label{sec:mechanisms}

The normalized displacement recurrence is the most efficient route to the
explicit coefficients in this article.  Two complementary formulations are
valuable both conceptually and computationally: a reduction to near-identity
reversion in the phase coordinate, and a perturbation expansion in a formal
coupling parameter.

\subsection{Phase-coordinate reduction and Lagrange--B\"urmann inversion}
\label{sec:phase-coordinate}

Let \(h=\Phi_0^{-1}\) denote the chosen Lambert-core inverse.  Introduce the
phase coordinate \(u=\Phi_0(x)\).  Then
\begin{equation}
 \Phi(h(u))=u+C+A(u),
 \qquad A(u)=\Psi(h(u)).
 \label{eq:phase-coordinate-map}
\end{equation}
After replacing \(Y\) by \(R=Y-C\), inversion of \(\Phi\) factors as
\begin{equation}
 \Phi^{-1}(Y)=h\!\left((I+A)^{-1}(R)\right),
 \qquad (I+A)(u)=u+A(u).
 \label{eq:inverse-factorization}
\end{equation}
Thus every Lambert-core problem can be reduced to the near-identity reversion
calculus already natural for polynomial-logarithmic transseries.

\begin{theorem}[Lagrange--B\"urmann phase reversion]
\label{thm:lagrange-phase}
Let \(\mathscr D\) be a differential transseries algebra in which the following
family is locally finite in the chosen valuation.  Then
\begin{equation}
 (I+A)^{-1}(R)
 =R+\sum_{n\ge1}\frac{(-1)^n}{n!}
 \left(\frac{\dd}{\dd R}\right)^{n-1}A(R)^n.
 \label{eq:lagrange-near-identity}
\end{equation}
Consequently,
\begin{equation}
 \Phi^{-1}(Y)
 =h\!\left(
 R+\sum_{n\ge1}\frac{(-1)^n}{n!}
 D_R^{n-1}A(R)^n
 \right).
 \label{eq:lagrange-core-composed}
\end{equation}
\end{theorem}

\begin{proof}
The classical Lagrange--B\"urmann identity for the inverse of
\(u\mapsto u+A(u)\) gives \cref{eq:lagrange-near-identity}.  In a formal
transseries algebra, the only additional issue is summability: at every
prescribed monomial, only finitely many summands may contribute.  This is the
stated local-finiteness hypothesis.  Composition with the already defined
core inverse \(h\) gives \cref{eq:lagrange-core-composed}.
\end{proof}

\begin{remark}
Formula \cref{eq:lagrange-near-identity} is universal but not always the best
coefficient engine.  Repeated differentiation expands expressions that the
normalized recurrence keeps factored.  It is most useful for structural
proofs, for deriving the first few perturbative terms, and for importing a
pre-existing near-identity transseries calculus.
\end{remark}

\subsection{Perturbation coefficients in closed differential form}

Insert a formal parameter \(\tau\) and solve
\begin{equation}
 \Phi_0(x)+\tau\Psi(x)=R,
 \qquad
 x=X+\tau x_1+\tau^2x_2+\tau^3x_3+\cdots.
 \label{eq:tau-perturbation}
\end{equation}
Write, at \(x=X\),
\[
 P=\Phi_0'(X),\quad Q=\Phi_0''(X),\quad S=\Phi_0'''(X),
 \qquad
 \Psi_j=\Psi^{(j)}(X).
\]
Direct Taylor expansion gives
\begin{align}
 x_1={}&-\frac{\Psi_0}{P},
 \label{eq:x1-general}\\
 x_2={}&\frac{\Psi_0\Psi_1}{P^2}
 -\frac{Q\Psi_0^2}{2P^3},
 \label{eq:x2-general}\\
 x_3={}&-\frac{\Psi_0\Psi_1^2}{P^3}
 -\frac{\Psi_0^2\Psi_2}{2P^3}
 +\frac{3Q\Psi_0^2\Psi_1}{2P^4}
 +\frac{S\Psi_0^3}{6P^4}
 -\frac{Q^2\Psi_0^3}{2P^5}.
 \label{eq:x3-general}
\end{align}
These identities are useful when \(\Psi\) is treated as one composite object.
The triangular \(q\)-recurrence is finer: it separates the different scales
inside \(\Psi\) and computes only the coefficient requested at each step.

\begin{proof}[Derivation]
Substitute \cref{eq:tau-perturbation} into the left-hand side and equate powers
of \(\tau\).  At first order one obtains \(Px_1+\Psi_0=0\).  At second order,
\[
 Px_2+\frac12Qx_1^2+\Psi_1x_1=0.
\]
At third order,
\[
 Px_3+Qx_1x_2+\frac16Sx_1^3+
 \Psi_1x_2+\frac12\Psi_2x_1^2=0.
\]
Solving successively and simplifying yields
\cref{eq:x1-general,eq:x2-general,eq:x3-general}.
\end{proof}

\subsection{Newton--Hensel lifting}

Define
\[
 F(\epsilon)=\Kcal_{p,v}(\epsilon)+\Rcal(q,v,\epsilon).
\]
Formal Newton iteration is
\begin{equation}
 \epsilon_{\mathrm{new}}
 =\epsilon-\frac{F(\epsilon)}{\partial_\epsilon F(\epsilon)}.
 \label{eq:newton-epsilon}
\end{equation}
Since \(\partial_\epsilon F(0)=v+1+O(q)\) is a unit, Newton lifting doubles the
known \(q\)-valuation under the usual complete-ring hypotheses.  A triangular
step is preferable when coefficients are wanted sequentially; Newton is
preferable for high-order truncated-series arithmetic because one inversion
of a unit replaces many scalar recurrences.

\section{Resonant logarithmic corrections}
\label{sec:resonance}

The large parameter \(v\) appears in two places: in the core kernel and in the
coefficients obtained by evaluating powers of \(\log x\).  Usually the latter
are lower order after division by the core slope.  A term containing exactly
one factor of \(\log x\), however, may survive coefficientwise as
\(v\to\infty\).  Barnes \(G\) exhibits this phenomenon.

\subsection{The coefficientwise limiting equation}

Suppose the normalized perturbation has an expansion
\begin{equation}
 \Rcal(q,v,E)=vH(q,E)+R_0(q,v,E),
 \qquad v^{-1}R_0\longrightarrow0
 \label{eq:resonant-decomposition}
\end{equation}
coefficientwise in the Hahn topology.  Divide
\cref{eq:normalized-equation} by \(v\).  Because
\begin{equation}
 \frac{1}{v}\Kcal_{p,v}(E)
 \longrightarrow \frac{(1+E)^p-1}{p},
 \label{eq:kernel-v-limit}
\end{equation}
the limiting correction \(E_0(q)\), when it exists, satisfies
\begin{equation}
 \frac{(1+E_0)^p-1}{p}+H(q,E_0)=0,
 \qquad E_0(0)=0.
 \label{eq:resonant-limit-equation}
\end{equation}

\begin{theorem}[Resonant coefficient limit]
\label{thm:resonant-limit}
Under the hypotheses of \cref{thm:triangular}, assume additionally
\cref{eq:resonant-decomposition} and suppose that the derivative of the
left-hand side of \cref{eq:resonant-limit-equation} at \(E_0=0,q=0\) is a
unit.  Then the coefficients of the unique solution \(\epsilon(q,v)\) tend,
as \(v\to\infty\), to the coefficients of the unique small solution
\(E_0(q)\) of \cref{eq:resonant-limit-equation}.
\end{theorem}

\begin{proof}
Apply the triangular recurrence coefficient by coefficient after division by
\(v\).  At each exponent, the coefficient equation has a limiting nonzero
linear divisor and depends only on already determined lower coefficients.
Induction therefore commutes with the coefficientwise limit.
\end{proof}

For a perturbation monomial
\(c x^{p-\rho}\log x\), substitution
\(x=X(1+E)\) gives the leading normalized contribution
\begin{equation}
 \frac{c}{ap}
 vq^\rho(1+E)^{p-\rho},
 \label{eq:resonant-monomial}
\end{equation}
so such monomials are precisely the elementary sources of resonance.  Terms
without \(\log x\) disappear from the limiting equation, while higher powers
of \(\log x\) require a modified scaling and may change the dominant core
rather than merely perturb it.

\section{Residuals, rigorous errors, and exponential sectors}
\label{sec:residuals}

\subsection{Residual-to-root transfer}

A symbolic transseries truncation becomes an analytic approximation only after
its residual is transferred through the inverse slope.

\begin{theorem}[Residual certificate]
\label{thm:residual-certificate}
Let \(\Phi\in C^2(I)\) be real and strictly monotone on an interval \(I\), let
\(x\in I\) solve \(\Phi(x)=Y\), and let \(\widehat x\in I\).  If
\(|\Phi'(t)|\ge m>0\) between \(x\) and \(\widehat x\), then
\begin{equation}
 |x-\widehat x|
 \le \frac{|\Phi(\widehat x)-Y|}{m}.
 \label{eq:mean-value-certificate}
\end{equation}
If the residual \(\delta=\Phi(\widehat x)-Y\) tends to zero and the relevant
derivatives have regular asymptotics, then
\begin{equation}
 x-\widehat x
 =-\frac{\delta}{\Phi'(\widehat x)}
 +O\!\left(
   \frac{\Phi''(\widehat x)\delta^2}{\Phi'(\widehat x)^3}
 \right).
 \label{eq:residual-newton-expansion}
\end{equation}
\end{theorem}

\begin{proof}
The first assertion is the mean-value theorem.  For the second, write
\(x=\widehat x+h\), Taylor-expand
\(0=\delta+\Phi'(\widehat x)h+\tfrac12\Phi''(\xi)h^2\), and use the first
bound to control \(h\).
\end{proof}

For \cref{eq:perturbed-phase}, provided \(\Psi'(X)=o(X^{p-1}v)\),
\begin{equation}
 \Phi'(X)\sim aX^{p-1}(v+1).
 \label{eq:asymptotic-inverse-slope}
\end{equation}
Thus a forward phase residual \(\delta\) produces an inverse error of leading
size
\begin{equation}
 -\frac{\delta}{aX^{p-1}(v+1)}.
 \label{eq:phase-error-division}
\end{equation}
This single division explains the \(1/\Lambda\) factor in the Gamma remainder
and the \(1/(U\lambda)\) response of the Barnes inverse to a phase error.

\subsection{Transport of exponentially small sectors}

Suppose an exponentially improved forward phase contains a sector
\begin{equation}
 E_\sigma(x)=S_\sigma(x)\exp(-\sigma x^\mu),
 \qquad \Re(\sigma x^\mu)>0.
 \label{eq:forward-exponential-sector}
\end{equation}
At first order, its inverse contribution is
\begin{equation}
 \Delta x_\sigma
 \sim-\frac{S_\sigma(X)\exp(-\sigma X^\mu)}
 {aX^{p-1}(v+1)}.
 \label{eq:inverse-exponential-sector}
\end{equation}
Products and substitutions of these sectors generate the deeper transseries
levels in the usual way.  Stokes multipliers are transported by the same
rule: evaluate the forward sector at the algebraic-logarithmic inverse and
divide by the inverse slope, then iterate if higher accuracy is required.

For Gamma and Barnes \(G\), explicit error bounds and exponentially improved
forward expansions are available in the literature
\cite{NemesGamma,NemesBarnes}.  The present article does not rederive their
terminant-function structure.  It identifies the exact algebraic operation
needed to convert any chosen forward hyperasymptotic truncation into an inverse
one.

\begin{warningbox}
An optimally truncated Bernoulli series may have an exponentially small
remainder, but the formal infinite Bernoulli series remains divergent.  The
exponential sector belongs after optimal truncation; it must not be appended
to an unqualified equality with the divergent series.
\end{warningbox}

\section{Complex asymptotic branches}
\label{sec:complex-branches}

The positive real formulas use \(\log y\) and \(\W_0\).  Complex inversion
requires two branch choices.  Let
\begin{equation}
 Y_m=\Log y+2\pi i m-C,
 \qquad m\in\mathbb Z,
 \label{eq:Y-complex-branch}
\end{equation}
where \(\Log\) is a fixed principal determination at the target.  For every
Lambert branch index \(k\), define
\begin{equation}
 v_{m,k}=\W_k\!\left(
  \frac{pY_m}{a}\e^{-p\beta}
 \right),
 \qquad
 X_{m,k}=\exp\!\left(\beta+\frac{v_{m,k}}{p}\right).
 \label{eq:complex-core-branches}
\end{equation}
The exponential form is preferable to an ambiguous \(p\)-th root: it keeps
the logarithm and Lambert branches compatible.  Whenever the forward phase
has a uniform asymptotic expansion in a sector containing \(X_{m,k}\), the
same normalized recurrence gives a formal inverse branch based at that core.
A Rouch\'e argument or a Newton contraction then identifies an actual local
inverse sheet for sufficiently large \(|Y_m|\).

For Gamma, the half-shifted variable and \(p=1\) leave the pair \((m,k)\) as
the asymptotic sheet labels.  For Barnes \(G\), \(p=2\) rotates the core by
approximately half the change in the Lambert phase.  Global branch topology is
more intricate because of poles or zeros of the original special function;
the construction here is an asymptotic local classification, not a claim that
the labels are globally one-to-one.

\section{The upper inverse Gamma transseries}
\label{sec:gamma}

\subsection{Why the half-shift is the correct normal form}

The standard Stirling expansion may be written, uniformly in sectors bounded
away from the negative real axis, as
\begin{equation}
 \log\Gamma(z+h)
 \sim\left(z+h-\frac12\right)\log z-z+\frac12\log(2\pi)
 +\sum_{k=2}^{\infty}
 \frac{(-1)^kB_k(h)}{k(k-1)z^{k-1}}.
 \label{eq:shifted-stirling-general}
\end{equation}
Set \(h=1/2\) and write
\begin{equation}
 t=x-\frac12.
 \label{eq:gamma-half-shift}
\end{equation}
Because \(B_{2k+1}(1/2)=0\), every even inverse power disappears and
\begin{equation}
 \log\Gamma\!\left(t+\frac12\right)
 \sim t(\log t-1)+\frac12\log(2\pi)
 +\sum_{k\ge1}\frac{a_k}{t^{2k-1}},
 \label{eq:gamma-shifted-phase}
\end{equation}
where
\begin{equation}
 a_k=\frac{B_{2k}(1/2)}{2k(2k-1)}
 =\frac{(2^{1-2k}-1)B_{2k}}{2k(2k-1)}.
 \label{eq:gamma-ak}
\end{equation}
The first values are
\begin{equation}
 a_1=-\frac1{24},\quad
 a_2=\frac7{2880},\quad
 a_3=-\frac{31}{40320},\quad
 a_4=\frac{127}{215040},\quad
 a_5=-\frac{511}{608256}.
 \label{eq:gamma-ak-first}
\end{equation}

The half-shift is more than a cosmetic acceleration.  It changes the support
of the normalized perturbation from a full integer grid to the even grid
\(2\N\).  Consequently, the displacement \(t-T\) contains only odd inverse
powers of the Lambert core.  This parity is exact to all algebraic orders.

\subsection{Core variables and normalized equation}

Let \(y\to+\infty\) and define
\begin{equation}
 R=\log y-\frac12\log(2\pi),
 \qquad
 w=\W_0(R/\e),
 \qquad
 T=\frac{R}{w}=\e^{w+1},
 \qquad
 \Lambda=\log T=w+1.
 \label{eq:gamma-core-full}
\end{equation}
Then
\begin{equation}
 T(\log T-1)=R,
 \qquad
 \frac{\dd}{\dd T}\bigl(T(\log T-1)\bigr)=\Lambda.
 \label{eq:gamma-core-identities}
\end{equation}
Put \(q=T^{-1}\) and
\begin{equation}
 t=T(1+\epsilon),
 \qquad
 \epsilon=\sum_{n\ge1}e_n(\Lambda)q^{2n}.
 \label{eq:gamma-epsilon}
\end{equation}
Subtract the core equation from \cref{eq:gamma-shifted-phase} and divide by
\(T\).  The exact formal equation is
\begin{equation}
 \Lambda\epsilon+
 (1+\epsilon)\log(1+\epsilon)-\epsilon
 +\sum_{k\ge1}a_kq^{2k}(1+\epsilon)^{-(2k-1)}=0.
 \label{eq:gamma-normalized}
\end{equation}
Since \(t-T=T\epsilon\), define
\begin{equation}
 d_{2n-1}(\Lambda)=e_n(\Lambda).
 \label{eq:gamma-d-e}
\end{equation}

\begin{theorem}[All-orders upper inverse Gamma expansion]
\label{thm:gamma-all-orders}
Let \(\Gamma_+^{-1}\) be the positive real inverse branch tending to infinity.
For every fixed \(N\ge0\),
\begin{equation}
 \Gamma_+^{-1}(y)
 =\frac12+T+\sum_{n=1}^{N}
 \frac{d_{2n-1}(\Lambda)}{T^{2n-1}}
 +O\!\left(\frac{1}{T^{2N+1}\Lambda}\right),
 \qquad y\to+\infty,
 \label{eq:gamma-all-orders}
\end{equation}
where \(T,\Lambda\) are given by \cref{eq:gamma-core-full}.  The coefficients
are uniquely determined by \(d_{2n-1}=e_n\) and
\begin{equation}
 e_n(\Lambda)=-\frac1\Lambda\coeff{q}{2n}
 \left[
 (1+E_{n-1})\log(1+E_{n-1})-E_{n-1}
 +\sum_{k=1}^{n}a_kq^{2k}
  (1+E_{n-1})^{-(2k-1)}
 \right],
 \label{eq:gamma-recurrence}
\end{equation}
with
\begin{equation}
 E_{n-1}=\sum_{j=1}^{n-1}e_j(\Lambda)q^{2j}.
 \label{eq:gamma-trunc-epsilon}
\end{equation}
Moreover,
\begin{equation}
 d_{2n-1}(\Lambda)
 \in \Lambda^{-1}\Q[\Lambda^{-1}],
 \qquad
 \deg_{\Lambda^{-1}}d_{2n-1}\le2n-1,
 \label{eq:gamma-coefficient-ring}
\end{equation}
and
\begin{equation}
 d_{2n-1}(\Lambda)=-\frac{a_n}{\Lambda}
 +O(\Lambda^{-2}).
 \label{eq:gamma-leading-block}
\end{equation}
\end{theorem}

\begin{proof}
\Cref{eq:gamma-normalized} is the specialization of
\cref{eq:normalized-equation} with \(p=a=\beta=1\) and support \(2\N\).
The coefficient of \(\epsilon\) is \(\Lambda\), so
\cref{thm:triangular} gives existence, uniqueness, and
\cref{eq:gamma-recurrence}.

We prove \cref{eq:gamma-coefficient-ring} by induction.  The nonlinear kernel
\((1+E)\log(1+E)-E\) starts with \(E^2/2\) and has rational coefficients.
The binomial expansions in the Bernoulli sector also have rational
coefficients.  If the lower \(e_j\) are polynomials in \(\Lambda^{-1}\) with
zero constant term and degree at most \(2j-1\), then every coefficient inside
the brackets in \cref{eq:gamma-recurrence} is a polynomial in
\(\Lambda^{-1}\) of degree at most \(2n-2\).  Division by \(\Lambda\) gives
the stated degree bound and zero constant term.  At exponent \(q^{2n}\), the
only contribution independent of lower \(e_j\) is \(a_nq^{2n}\), which proves
\cref{eq:gamma-leading-block}.

For the analytic remainder, use the duplication identity
\[
 \log\Gamma\!\left(t+\frac12\right)
 =\log\Gamma(2t)-\log\Gamma(t)+(1-2t)\log2+\frac12\log\pi.
\]
Apply the standard remainder bound separately to the two ordinary
positive-real Stirling expansions.  For every fixed \(M\), this shows that
\cref{eq:gamma-shifted-phase} truncated after \(a_Mt^{-(2M-1)}\) has
remainder \(O(t^{-(2M+1)})\).  Choose \(M=N+1\).  Substitution of the
formal inverse through \(e_Nq^{2N}\) then leaves a forward phase residual
\(O(T^{-(2N+1)})\).  The inverse phase slope is
\(\Lambda+O(T^{-2})\).  Applying \cref{thm:residual-certificate} divides the
phase residual by \(\Lambda\), yielding the error in
\cref{eq:gamma-all-orders}.
\end{proof}

\subsection{Explicit correction blocks}

The first five nonzero blocks are
\begin{align}
 d_1={}&\frac{1}{24\Lambda},
 \label{eq:gamma-d1}\\[0.3em]
 d_3={}&-\frac{14\Lambda^2+10\Lambda+5}{5760\Lambda^3},
 \label{eq:gamma-d3}\\[0.3em]
 d_5={}&\frac{
 2232\Lambda^4+1176\Lambda^3+714\Lambda^2+350\Lambda+105}
 {2903040\Lambda^5},
 \label{eq:gamma-d5}\\[0.3em]
 d_7={}&-\frac{1}{1393459200\Lambda^7}
 \bigl(
 822960\Lambda^6+292536\Lambda^5+136956\Lambda^4
 \notag\\[-0.2em]
 &\hspace{4.5em}{}+68040\Lambda^3+32970\Lambda^2
 +12250\Lambda+2625
 \bigr),
 \label{eq:gamma-d7}\\[0.3em]
 d_9={}&\frac{1}{367873228800\Lambda^9}
 \bigl(
 309052800\Lambda^8+77920128\Lambda^7+26024592\Lambda^6
 \notag\\[-0.2em]
 &\hspace{2.5em}{}+10019592\Lambda^5+4516028\Lambda^4
 +2023560\Lambda^3+812350\Lambda^2
 \notag\\[-0.2em]
 &\hspace{2.5em}{}+242550\Lambda+40425
 \bigr).
 \label{eq:gamma-d9}
\end{align}
Thus, through the first three correction blocks,
\begin{align}
 \Gamma_+^{-1}(y)={}&\frac12+T
 +\frac{1}{24T\Lambda}
 -\frac{14\Lambda^2+10\Lambda+5}{5760T^3\Lambda^3}
 \notag\\
 &+\frac{2232\Lambda^4+1176\Lambda^3+714\Lambda^2
 +350\Lambda+105}{2903040T^5\Lambda^5}
 +O\!\left(\frac{1}{T^7\Lambda}\right).
 \label{eq:gamma-three-blocks}
\end{align}

For comparison with expansions written as separate powers of \(\log T\), the
first blocks are
\begin{align}
 d_3={}&-\frac{7}{2880\Lambda}
        -\frac{1}{576\Lambda^2}
        -\frac{1}{1152\Lambda^3},
 \label{eq:gamma-d3-expanded}\\
 d_5={}&\frac{31}{40320\Lambda}
 +\frac{7}{17280\Lambda^2}
 +\frac{17}{69120\Lambda^3}
 +\frac{5}{41472\Lambda^4}
 +\frac{1}{27648\Lambda^5}.
 \label{eq:gamma-d5-expanded}
\end{align}
The expansion of Gerhold, Hubalek, and Tomovski is precisely
\cref{eq:gamma-three-blocks} through \(d_3/T^3\), expressed in the separated
form \cref{eq:gamma-d3-expanded} \cite{GerholdHubalekTomovski}.

\subsection{A hand derivation of the first two blocks}

The recurrence is transparent at low order.  Since
\[
 (1+E)\log(1+E)-E=\frac12E^2-\frac16E^3+\cdots,
\]
the coefficient of \(q^2\) in \cref{eq:gamma-normalized} is
\(\Lambda e_1+a_1\).  Therefore
\[
 e_1=-\frac{a_1}{\Lambda}=\frac{1}{24\Lambda}.
\]
At order \(q^4\), the contributions are
\[
 \Lambda e_2+\frac12e_1^2-a_1e_1+a_2=0.
\]
Substituting \(a_1=-1/24\), \(a_2=7/2880\), and
\(e_1=1/(24\Lambda)\) yields \cref{eq:gamma-d3}.  Higher orders are no more
conceptually difficult, but the number of partitions of the target exponent
grows; truncated-series arithmetic is preferable to hand expansion.

\subsection{Outer expansion in iterated logarithms}
\label{sec:gamma-outer}

The Lambert core itself can be expanded without first writing the generic
\(\W\)-series.  Let
\begin{equation}
 L=\log R,
 \qquad \ell=\log L.
 \label{eq:gamma-L-ell}
\end{equation}
Writing \(T=(R/L)S\) in \(T(\log T-1)=R\) gives the implicit equation
\begin{equation}
 S\left(1-\frac{\ell+1}{L}+\frac{\log S}{L}\right)=1.
 \label{eq:gamma-S-equation}
\end{equation}
Ordinary formal reversion in \(L^{-1}\) yields
\begin{align}
 T=\frac{R}{L}\biggl[{}&1+\frac{\ell+1}{L}
 +\frac{\ell(\ell+1)}{L^2}
 +\frac{(\ell-1)(\ell+1)(2\ell+1)}{2L^3}
 \notag\\
 &+\frac{(\ell+1)(6\ell^3-8\ell^2-7\ell+1)}{6L^4}
 \notag\\
 &+\frac{(\ell+1)(12\ell^4-29\ell^3-17\ell^2+17\ell+5)}
 {12L^5}
 +O\!\left(\frac{\ell^6}{L^6}\right)
 \biggr].
 \label{eq:gamma-T-outer}
\end{align}
Combining \cref{eq:gamma-T-outer} with \cref{eq:gamma-all-orders} gives a
pure iterated-logarithmic expansion if desired.  The ordering must be read
lexicographically:
\begin{equation}
 \frac{T}{L^N}\longrightarrow\infty
 \quad\text{for every fixed }N,
 \qquad
 1\gg \frac1{T\Lambda}.
 \label{eq:gamma-hierarchy}
\end{equation}
Hence the shift \(1/2\) is beyond all orders of the outer \(L^{-1}\) scale,
and the Bernoulli inverse blocks are beyond that shift.

\subsection{Numerical verification}
\label{sec:gamma-numerics}

To avoid numerical inversion error, the following test starts from an exact
chosen abscissa \(x\), forms \(y=\Gamma(x)\), reconstructs \(T,\Lambda\) from
\(y\), and compares successive truncations with the original \(x\).  Let
\(E_0\) denote approximation minus exact value for \(1/2+T\), and let
\(E_j\) use the same sign convention after including the first \(j\)
correction blocks \(d_1/T,d_3/T^3,\ldots\).

\begin{table}[htbp]
\centering
\small
\caption{Signed errors (approximation minus exact value) for the upper inverse Gamma expansion.}
\label{tab:gamma-errors}
\begin{tabular}{r@{\quad}rrrrr}
\toprule
\(x\) & \(E_0\) & \(E_1\) & \(E_2\) & \(E_3\) & \(E_4\)\\
\midrule
5   & \(-6.1414\!\times10^{-3}\) & \( 2.8694\!\times10^{-5}\)
    & \(-4.1862\!\times10^{-7}\) & \( 1.3452\!\times10^{-8}\)
    & \(-8.2694\!\times10^{-10}\)\\
10  & \(-1.9470\!\times10^{-3}\) & \( 1.7430\!\times10^{-6}\)
    & \(-5.7560\!\times10^{-9}\) & \( 4.4485\!\times10^{-11}\)
    & \(-6.5578\!\times10^{-13}\)\\
20  & \(-7.1924\!\times10^{-4}\) & \( 1.4126\!\times10^{-7}\)
    & \(-1.1183\!\times10^{-10}\) & \( 2.1113\!\times10^{-13}\)
    & \(-7.5612\!\times10^{-16}\)\\
100 & \(-9.1031\!\times10^{-5}\) & \( 6.2869\!\times10^{-10}\)
    & \(-1.9388\!\times10^{-14}\) & \( 1.4440\!\times10^{-18}\)
    & \(-2.0231\!\times10^{-22}\)\\
\bottomrule
\end{tabular}
\end{table}

The clean reduction by several orders per block is typical before optimal
truncation.  At a fixed moderate \(x\), sufficiently remote Bernoulli blocks
will eventually grow because the underlying Stirling series is divergent.

\section{The other positive inverse Gamma branch}
\label{sec:gamma-small-branch}

For completeness, a large positive \(y\) also has a positive preimage tending
to zero.  This branch is governed by a local pole, not by the Lambert core.
Put \(s=1/y\).  Since
\begin{equation}
 s=\frac{x}{\Gamma(1+x)},
 \label{eq:small-gamma-relation}
\end{equation}
Lagrange inversion gives the convergent local series
\begin{equation}
 \Gamma_0^{-1}(y)
 =\sum_{n\ge1}\frac{s^n}{n!}
 \left.
 \frac{\dd^{n-1}}{\dd x^{n-1}}\Gamma(1+x)^n
 \right|_{x=0},
 \qquad s=\frac1y.
 \label{eq:gamma-small-lagrange}
\end{equation}
Using
\begin{equation}
 \log\Gamma(1+x)
 =-\gamma x+\sum_{m\ge2}\frac{(-1)^m\zeta(m)}{m}x^m,
 \qquad |x|<1,
 \label{eq:loggamma-near-one}
\end{equation}
we find
\begin{align}
 \Gamma_0^{-1}(y)={}&\frac1y-\frac{\gamma}{y^2}
 +\left(\frac{3\gamma^2}{2}+\frac{\pi^2}{12}\right)\frac1{y^3}
 \notag\\
 &-\left(\frac{8\gamma^3}{3}+\frac{\gamma\pi^2}{3}
 +\frac{\zeta(3)}{3}\right)\frac1{y^4}
 \notag\\
 &+\left(\frac{125\gamma^4}{24}
 +\frac{25\gamma^2\pi^2}{24}
 +\frac{5\gamma\zeta(3)}{3}
 +\frac{29\pi^4}{1440}\right)\frac1{y^5}
 +O(y^{-6}).
 \label{eq:gamma-small-explicit}
\end{align}
More generally, every pole at \(-n\) gives a local inverse branch with leading
term
\begin{equation}
 x=-n+\frac{(-1)^n}{n!\,y}+O(y^{-2}).
 \label{eq:gamma-pole-branches}
\end{equation}
These pole-centered branches belong to ordinary local series reversion and are
algebraically separate from the large branch analyzed in
\cref{sec:gamma}.

\section{The inverse Barnes \texorpdfstring{\(G\)}{G}-transseries}
\label{sec:barnes}

\subsection{Forward phase in a centered variable}

Barnes' \(G\)-function is characterized by
\begin{equation}
 G(z+1)=\Gamma(z)G(z),
 \qquad G(1)=1.
 \label{eq:barnes-recurrence-definition}
\end{equation}
For inversion on the large positive real axis, write the argument as \(z+1\)
from the outset.  The classical forward expansion is
\begin{align}
 \log G(z+1)\sim{}&
 \frac12z^2\log z-\frac34z^2
 +\frac12\log(2\pi)\,z
 -\frac1{12}\log z
 +\zeta'(-1)
 \notag\\
 &+\sum_{k\ge1}\frac{B_{2k+2}}{4k(k+1)z^{2k}},
 \qquad z\to\infty,
 \label{eq:barnes-forward}
\end{align}
valid in every fixed sector avoiding the negative real axis
\cite{DLMF517,NemesBarnes}.  Define
\begin{equation}
 \alpha=\frac12\log(2\pi),
 \qquad b=-\frac1{12},
 \qquad C_G=\zeta'(-1)=\frac1{12}-\log A,
 \label{eq:barnes-alpha-b-C}
\end{equation}
where \(A\) is the Glaisher--Kinkelin constant, and
\begin{equation}
 c_k=\frac{B_{2k+2}}{4k(k+1)}.
 \label{eq:barnes-ck}
\end{equation}
Then
\begin{equation}
 c_1=-\frac1{240},\qquad
 c_2=\frac1{1008},\qquad
 c_3=-\frac1{1440},\qquad
 c_4=\frac1{1056}.
 \label{eq:barnes-ck-first}
\end{equation}
In the notation of the general model,
\begin{equation}
 a=\frac12,
 \qquad p=2,
 \qquad \beta=\frac32,
 \label{eq:barnes-core-parameters}
\end{equation}
and the perturbation is
\begin{equation}
 \Psi(z)=\alpha z+b\log z+\sum_{k\ge1}c_kz^{-2k}.
 \label{eq:barnes-Psi}
\end{equation}

\subsection{Exact core variables}

Let \(G(z+1)=y\) and set
\begin{equation}
 R=\log y-C_G.
 \label{eq:barnes-R}
\end{equation}
The core equation
\begin{equation}
 \frac12U^2\left(\log U-\frac32\right)=R
 \label{eq:barnes-core-equation}
\end{equation}
is solved exactly by
\begin{equation}
 v=\W_0(4R/\e^3),
 \qquad
 U=2\sqrt{\frac{R}{v}}=\e^{(v+3)/2},
 \qquad
 \lambda=\log U-1=\frac{v+1}{2}.
 \label{eq:barnes-U-v-lambda}
\end{equation}
The core slope is
\begin{equation}
 \frac{\dd}{\dd U}
 \left[\frac12U^2\left(\log U-\frac32\right)\right]
 =U(\log U-1)=U\lambda.
 \label{eq:barnes-core-slope}
\end{equation}

Put \(q=U^{-1}\) and
\begin{equation}
 z=U(1+\epsilon),
 \qquad
 \epsilon=\sum_{n\ge1}e_n(\lambda)q^n.
 \label{eq:barnes-epsilon}
\end{equation}
After subtraction of \cref{eq:barnes-core-equation} and division by
\(U^2/2\), the exact normalized equation is
\begin{align}
 0={}&\Kcal_\lambda(\epsilon)
 +2\alpha q(1+\epsilon)
 +2bq^2\bigl(\lambda+1+\log(1+\epsilon)\bigr)
 \notag\\
 &+2\sum_{k\ge1}c_kq^{2k+2}(1+\epsilon)^{-2k},
 \label{eq:barnes-normalized}
\end{align}
where
\begin{equation}
 \Kcal_\lambda(E)
 =(1+E)^2\left(\lambda-\frac12+\log(1+E)\right)
 -\left(\lambda-\frac12\right).
 \label{eq:barnes-kernel}
\end{equation}
Its linear part is
\begin{equation}
 \Kcal_\lambda(E)=2\lambda E+O(E^2).
 \label{eq:barnes-kernel-linear}
\end{equation}
Since
\[
 z=U+e_1+e_2U^{-1}+e_3U^{-2}+\cdots,
\]
we set
\begin{equation}
 b_n(\lambda)=e_{n+1}(\lambda),
 \qquad n\ge0.
 \label{eq:barnes-b-e}
\end{equation}

\begin{theorem}[All-orders inverse Barnes expansion]
\label{thm:barnes-all-orders}
Let \(\Bcal(y)\) be the eventual positive solution of \(G(z+1)=y\).  For every
fixed \(N\ge0\),
\begin{equation}
 \Bcal(y)=U+\sum_{n=0}^{N}\frac{b_n(\lambda)}{U^n}
 +O(U^{-N-1}),
 \qquad y\to+\infty,
 \label{eq:barnes-all-orders}
\end{equation}
where \(U,\lambda\) are given by \cref{eq:barnes-U-v-lambda}.  Define
\begin{equation}
 E_{n-1}=\sum_{j=1}^{n-1}e_j(\lambda)q^j.
 \label{eq:barnes-E-trunc}
\end{equation}
Then \(e_n\), and hence every \(b_n\), is uniquely determined by
\begin{align}
 e_n(\lambda)=-\frac1{2\lambda}\coeff{q}{n}
 \biggl[{}&\Kcal_\lambda(E_{n-1})-2\lambda E_{n-1}
 +2\alpha q(1+E_{n-1})
 \notag\\
 &+2bq^2\bigl(\lambda+1+\log(1+E_{n-1})\bigr)
 \notag\\
 &+2\sum_{1\le k\le(n-2)/2}
 c_kq^{2k+2}(1+E_{n-1})^{-2k}
 \biggr].
 \label{eq:barnes-recurrence}
\end{align}
The last sum is empty when \(n<4\).  The coefficients satisfy
\begin{equation}
 e_n\in\Q[\alpha,\lambda,\lambda^{-1}],
 \qquad
 e_n(-\alpha)=(-1)^ne_n(\alpha),
 \label{eq:barnes-ring-parity}
\end{equation}
and are bounded coefficient functions as \(\lambda\to+\infty\).
\end{theorem}

\begin{proof}
\Cref{eq:barnes-normalized} has support \(\N\) and linear divisor
\(2\lambda\).  The general triangular theorem gives
\cref{eq:barnes-recurrence}.  The coefficient ring follows by induction from
the rational Taylor coefficients of the logarithm and binomial powers.  The
transformation
\begin{equation}
 (q,\alpha,E)\longmapsto(-q,-\alpha,E)
 \label{eq:barnes-parity-transform}
\end{equation}
leaves \cref{eq:barnes-normalized} invariant.  Uniqueness therefore implies
\(\epsilon(-q,-\alpha)=\epsilon(q,\alpha)\), which is equivalent to the
parity assertion.

To prove boundedness, divide \cref{eq:barnes-normalized} by \(\lambda\).
The only coefficient growing linearly with \(\lambda\), besides the core
kernel, is the constant part of the \(bq^2\)-term.  The limiting equation has a
small formal solution with bounded coefficients; \cref{thm:resonant-limit}
then gives coefficientwise convergence.  The forward Barnes remainder can be
truncated to any prescribed algebraic order.  Its transfer through the slope
\(U\lambda(1+o(1))\), together with boundedness of the next coefficient block,
gives the error in \cref{eq:barnes-all-orders}.
\end{proof}

\subsection{Explicit Barnes coefficient blocks}

The first coefficients are
\begin{align}
 b_0={}&-\frac{\alpha}{\lambda},
 \label{eq:barnes-b0}\\[0.4em]
 b_1={}&\frac{6\alpha^2\lambda-6\alpha^2+\lambda^3+\lambda^2}
 {12\lambda^3},
 \label{eq:barnes-b1}\\[0.4em]
 b_2={}&\frac{\alpha(8\alpha^2\lambda-6\alpha^2+\lambda^2)}
 {12\lambda^5}.
 \label{eq:barnes-b2}
\end{align}
For the next two, define
\begin{align}
 P_3(\lambda,\alpha)={}&
 180\alpha^4\lambda^3+120\alpha^4\lambda^2
 -1500\alpha^4\lambda+900\alpha^4
 \notag\\
 &+60\alpha^2\lambda^5+120\alpha^2\lambda^4
 +60\alpha^2\lambda^3-180\alpha^2\lambda^2
 \notag\\
 &+5\lambda^7-\lambda^6+5\lambda^5+5\lambda^4,
 \label{eq:barnes-P3}\\[0.4em]
 P_4(\lambda,\alpha)={}&
 576\alpha^4\lambda^3+480\alpha^4\lambda^2
 -2520\alpha^4\lambda+1260\alpha^4
 \notag\\
 &+80\alpha^2\lambda^5+280\alpha^2\lambda^4
 +200\alpha^2\lambda^3-300\alpha^2\lambda^2
 \notag\\
 &-22\lambda^7-11\lambda^6+10\lambda^5+15\lambda^4.
 \label{eq:barnes-P4}
\end{align}
Then
\begin{equation}
 b_3=-\frac{P_3(\lambda,\alpha)}{1440\lambda^7},
 \qquad
 b_4=-\frac{\alpha P_4(\lambda,\alpha)}{1440\lambda^9}.
 \label{eq:barnes-b3-b4}
\end{equation}
The longer coefficient \(b_5\) is listed in
\cref{app:barnes-b5}.  The expansion through \(U^{-2}\) is therefore
\begin{align}
 \Bcal(y)={}&U-\frac{\alpha}{\lambda}
 +\frac{6\alpha^2\lambda-6\alpha^2+\lambda^3+\lambda^2}
 {12\lambda^3U}
 \notag\\
 &+\frac{\alpha(8\alpha^2\lambda-6\alpha^2+\lambda^2)}
 {12\lambda^5U^2}
 +O(U^{-3}).
 \label{eq:barnes-three-blocks}
\end{align}

\subsection{Resonance and the limiting algebraic subseries}
\label{sec:barnes-resonance}

The resonance is completely explicit.  Divide
\cref{eq:barnes-normalized} by \(2\lambda\) and let
\(\lambda\to+\infty\) coefficientwise.  The \(\alpha q\)-term, all
\(c_k\)-terms, and the nonconstant part of the logarithmic term vanish.  Since
\[
 \frac{\Kcal_\lambda(E)}{2\lambda}
 \longrightarrow\frac{(1+E)^2-1}{2},
\]
the limiting series \(E_0(q)\) satisfies
\begin{equation}
 \frac{(1+E_0)^2-1}{2}+bq^2=0.
 \label{eq:barnes-limit-equation}
\end{equation}
The branch with \(E_0(0)=0\) is
\begin{equation}
 E_0(q)=\sqrt{1-2bq^2}-1
 =\sqrt{1+\frac{q^2}{6}}-1.
 \label{eq:barnes-limit-series}
\end{equation}
Therefore
\begin{equation}
 \lim_{\lambda\to\infty}e_{2m+1}(\lambda)=0,
 \qquad
 \lim_{\lambda\to\infty}e_{2m}(\lambda)
 =\binom{1/2}{m}6^{-m}.
 \label{eq:barnes-e-limits}
\end{equation}
Equivalently,
\begin{equation}
 \lim_{\lambda\to\infty}b_{2m}(\lambda)=0,
 \qquad
 \lim_{\lambda\to\infty}b_{2m-1}(\lambda)
 =\binom{1/2}{m}6^{-m}.
 \label{eq:barnes-b-limits}
\end{equation}
The first limiting constants are
\begin{equation}
 b_1\to\frac1{12},
 \qquad
 b_3\to-\frac1{288},
 \qquad
 b_5\to\frac1{3456}.
 \label{eq:barnes-limits-first}
\end{equation}
This coefficientwise algebraic limit is invisible if the forward
\(-\tfrac1{12}\log z\) term is treated as an undifferentiated small error.
It is also a useful check on symbolic computations: every odd inverse-power
block must approach the corresponding binomial coefficient, while every even
block must vanish.

\subsection{Outer iterated-logarithmic expansion}
\label{sec:barnes-outer}

Let
\begin{equation}
 Z=\frac{4R}{\e^3},
 \qquad
 L_1=\log Z,
 \qquad
 L_2=\log L_1.
 \label{eq:barnes-L1-L2}
\end{equation}
Expanding \(v=\W_0(Z)\) and then \(v^{-1/2}\) gives
\begin{align}
 U=2\sqrt{\frac{R}{L_1}}\biggl[{}&1+\frac{L_2}{2L_1}
 +\frac{L_2(3L_2-4)}{8L_1^2}
 +\frac{L_2(5L_2^2-16L_2+8)}{16L_1^3}
 \notag\\
 &+\frac{L_2(105L_2^3-568L_2^2+720L_2-192)}{384L_1^4}
 +O\!\left(\frac{L_2^5}{L_1^5}\right)
 \biggr].
 \label{eq:barnes-U-outer}
\end{align}
For the inverse of \(G\) itself, the nested hierarchy begins
\begin{equation}
 G_+^{-1}(y)
 =U+1-\frac{\alpha}{\lambda}
 +\frac{b_1(\lambda)}{U}+\cdots.
 \label{eq:barnes-full-hierarchy}
\end{equation}
For every fixed \(N\), \(U/L_1^N\to\infty\); hence the shift \(+1\) is
beyond all orders of the outer iterated-logarithmic scale.  The displacement
\(-\alpha/\lambda\) is smaller than the shift but larger than every inverse
power of \(U\).  This four-level ordering is worth displaying explicitly:
\begin{equation}
 \frac{U}{L_1^N}\gg1\gg\lambda^{-1}\gg U^{-1}
 \qquad(N\text{ fixed}).
 \label{eq:barnes-hierarchy}
\end{equation}

\subsection{Numerical verification}
\label{sec:barnes-numerics}

Choose an exact \(z\), form \(y=G(z+1)\), reconstruct \(U,\lambda\), and let
\(E_\star\) be core approximation minus exact value.  Let \(E_j\) use the
same sign convention after including \(b_0,b_1/U,\ldots,b_j/U^j\).

\begin{table}[htbp]
\centering
\small
\caption{Signed errors (approximation minus exact value) for \(\Bcal(y)\), where \(y=G(z+1)\).}
\label{tab:barnes-errors}
\resizebox{\textwidth}{!}{%
\begin{tabular}{r@{\quad}rrrrrr}
\toprule
\(z\) & \(E_\star\) & \(E_0\) & \(E_1\) & \(E_2\) & \(E_3\) & \(E_4\)\\
\midrule
5 & \( 1.1187\) & \(-1.3867\!\times10^{-2}\)
  & \(-7.8343\!\times10^{-3}\) & \(-1.5935\!\times10^{-3}\)
  & \(-1.8345\!\times10^{-4}\) & \( 5.8515\!\times10^{-6}\)\\
10 & \( 6.5273\!\times10^{-1}\) & \(-2.0081\!\times10^{-2}\)
   & \(-8.4026\!\times10^{-4}\) & \( 1.5249\!\times10^{-5}\)
   & \( 3.2299\!\times10^{-6}\) & \( 4.1936\!\times10^{-8}\)\\
20 & \( 4.4669\!\times10^{-1}\) & \(-8.7205\!\times10^{-3}\)
   & \(-6.6821\!\times10^{-5}\) & \( 2.3736\!\times10^{-6}\)
   & \( 3.4133\!\times10^{-8}\) & \(-1.5564\!\times10^{-9}\)\\
100 & \( 2.5342\!\times10^{-1}\) & \(-1.2959\!\times10^{-3}\)
    & \(-3.9481\!\times10^{-7}\) & \( 8.1326\!\times10^{-9}\)
    & \(-5.3660\!\times10^{-12}\) & \(-8.9993\!\times10^{-14}\)\\
\bottomrule
\end{tabular}%
}
\end{table}

The core error is much larger than in the Gamma example because the linear
term \(\alpha z\) produces the slowly decaying displacement
\(-\alpha/\lambda\).  Once that block is included, the triangular expansion
again improves rapidly.  At \(z=5\), adding still more terms eventually
encounters the asymptotic nature of the forward Bernoulli sector; at larger
arguments the useful range grows.

\section{General phase families of Gamma--Barnes type}
\label{sec:general-families}

\subsection{Polynomial-logarithmic perturbations below the leading power}

A broad class covered directly by \cref{thm:triangular} is
\begin{equation}
 \Phi(x)=a x^p(\log x-\beta)+C
 +\sum_{\rho\in A}x^{p-\rho}P_\rho(\log x)
 +\text{smaller sectors},
 \label{eq:general-polylog-phase}
\end{equation}
where \(A\subset\R_{>0}\) is well ordered and each \(P_\rho\) is a
polynomial or an asymptotic coefficient transseries.  Substitution
\(x=X(1+E)\), with \(q=X^{-1}\), gives
\begin{align}
 \Rcal(q,v,E)={}&\frac1a\sum_{\rho\in A}q^\rho
 (1+E)^{p-\rho}
 P_\rho\!\left(\beta+\frac vp+\log(1+E)\right)
 \notag\\
 &+\text{smaller normalized sectors}.
 \label{eq:general-polylog-normalized}
\end{align}
Every term has positive \(q\)-valuation.  Powers of \(v\) in its coefficient
are harmless because \(q^\rho v^M\to0\) for every fixed \(M\).  The inverse
support lies in the additive semigroup generated by \(A\).

This immediately explains the two examples:
\begin{center}
\begin{tabular}{lllll}
\toprule
Function & shifted variable & \(p\) & base support & first perturbation\\
\midrule
\(\Gamma(t+1/2)\) & \(t=x-1/2\) & 1 & \(2\N\) & \(-1/(24t)\)\\
\(G(z+1)\) & \(z=x-1\) & 2 & \(\N\) & \(\alpha z\)\\
\bottomrule
\end{tabular}
\end{center}
For Gamma, a drop of two units in normalized power occurs because the forward
term \(t^{-1}\) is divided by the order-\(t\) core.  For Barnes, the forward
linear term \(\alpha z\) is divided by the order-\(z^2\) core and produces
\(q^1\), so the full integer grid is generated.

\subsection{What must be absorbed into the core}

Terms with the same algebraic power \(x^p\) as the leading phase are not
\(q\)-small.  For example,
\[
 a x^p\log x+c x^p=a x^p(\log x-\beta),
 \qquad \beta=-c/a,
\]
so the \(cx^p\) term should be absorbed exactly into \(\beta\).  More generally,
a phase
\begin{equation}
 x^p\left(a_r(\log x)^r+a_{r-1}(\log x)^{r-1}+\cdots\right)
 \label{eq:same-power-log-poly}
\end{equation}
requires an outer inversion in the logarithmic coefficient scale before the
\(q\)-grid theorem is applied.  Treating a same-power term as a small
perturbation makes \(\Rcal(q,v,0)\) have valuation zero and violates the
hypothesis of \cref{thm:triangular}.

This gives a practical normalization rule:
\begin{enumerate}
\item collect and invert the entire dominant algebraic power, including the
      necessary logarithmic polynomial;
\item use drops in algebraic power to define the exponentially finer \(q\)-grid;
\item place optimally truncated or resurgent exponentials in still deeper
      sectors.
\end{enumerate}

\subsection{Exact Lambert inversion for a logarithmic-power core}
\label{sec:log-power-core}

The exact core need not be linear in \(\log x\).  Let
\begin{equation}
 \Phi_{0,r}(x)=a x^p(\log x-\beta)^r,
 \qquad p\ne0,\quad r\ne0.
 \label{eq:log-power-core}
\end{equation}
Choose compatible branches of powers and logarithms.  For
\(\Phi_{0,r}(X)=R\), define
\begin{equation}
 Z=\frac pr\left(\frac{R}{a\e^{p\beta}}\right)^{1/r},
 \qquad
 w=\W_k(Z),
 \qquad
 u=\frac rp w,
 \qquad
 X=\e^{\beta+u}.
 \label{eq:log-power-core-inverse}
\end{equation}
Then \(\Phi_{0,r}(X)=R\).  Indeed,
\[
 u^r\e^{pu}
 =\left(\frac rp\right)^rw^r\e^{rw}
 =\frac{R}{a\e^{p\beta}}.
\]
The derivative is
\begin{equation}
 \Phi_{0,r}'(X)
 =aX^{p-1}u^{r-1}(pu+r).
 \label{eq:log-power-core-slope}
\end{equation}
With \(x=X(1+E)\), the normalized universal kernel becomes
\begin{equation}
 \Kcal_{p,r,u}(E)
 =(1+E)^p
 \left(1+\frac{\log(1+E)}{u}\right)^r-1,
 \label{eq:log-power-kernel}
\end{equation}
whose linear coefficient is
\begin{equation}
 p+\frac ru=\frac{pu+r}{u}.
 \label{eq:log-power-linear}
\end{equation}
The same Hahn implicit-function proof applies whenever this coefficient is a
unit and the normalized perturbation has positive valuation.  The case
\(r=1\) reduces, after multiplication by \(v/p\), to the kernel used in the
rest of the article.

This extension covers inverse phases led by
\(x^p(\log x)^r\) without introducing a new special inverse function.  It also
clarifies the boundary of the method: a genuinely non-polynomial function of
\(\log x\) may still be invertible in a transseries coefficient field, but the
Lambert closed form can disappear.

\subsection{Coefficient algebras and nested supports}

For computational work it is useful to distinguish three coefficient choices.

\paragraph{Rational-function coefficients.}
Keeping \(v\) or \(\lambda\) exact gives coefficients in fields such as
\(\Q(v)\) or \(\Q(\alpha,\lambda)\).  This is the compact form used for the
explicit Gamma and Barnes blocks.

\paragraph{Laurent coefficients at large slope.}
Expanding each rational coefficient in \(v^{-1}\) gives an element of
\(\Q((v^{-1}))[[q^S]]\).  The order is lexicographic: compare \(q\)-exponents
first, then inverse powers of \(v\).  This form exposes the leading Newton
response and resonant limits.

\paragraph{Full outer logarithmic coefficients.}
Substituting the expansion of \(v=\W(Z)\) embeds the result into a larger
polynomial-logarithmic transseries in \(L_1^{-1}\), \(L_2=\log L_1\), and
\(q\).  Since \(q\) itself is an exponential of \(-v/p\), the resulting
monomial group has at least two logarithmic-exponential levels.  Expansion in
this largest algebra should be postponed until presentation or comparison
requires it; it is usually inefficient for coefficient generation.

\subsection{Higher Barnes multiple gamma functions}

The logarithms of higher Barnes multiple gamma functions have leading phases
of the schematic form
\begin{equation}
 a_m x^m\log x+b_mx^m+\text{lower algebraic powers and logarithms}.
 \label{eq:multiple-gamma-schematic}
\end{equation}
After absorbing \(b_mx^m\) into a shift of the logarithm, the core has
\(p=m\) and \(r=1\).  Therefore the same Lambert-core and Hahn-grid machinery
applies, subject to the sign and convention used for the particular multiple
gamma normalization.  The support semigroup is read directly from the drops
in algebraic degree in the known forward expansion.  Barnes \(G\) is the
first nontrivial member, with \(m=2\).

\section{Coefficient computation and numerical use}
\label{sec:algorithms}

\subsection{A generic truncated-series algorithm}

The recurrence \cref{eq:general-triangular-recurrence} is readily implemented
in any computer algebra system supporting truncated Hahn or Puiseux series.
For an integral grid, the following pseudocode is sufficient.

\begin{algorithm}[Triangular inversion on a \(q\)-grid]
\label{alg:generic-triangular}
Given a truncation order \(N\), a kernel
\(K(E)=(v+1)E+K_{\ge2}(E)\), and a normalized perturbation
\(R(q,E)\):
\begin{lstlisting}
E := 0
for n := 1,...,N:
    H := truncate_q(K_ge_2(E) + R(q,E), order = n)
    e_n := - coefficient(H, q^n) / (v+1)
    E := E + e_n q^n
return E
\end{lstlisting}
At step \(n\), every operation may be truncated immediately above \(q^n\).
\end{algorithm}

If the base support is \(2\N\), as for shifted Gamma, replace \(q^n\) by
\(q^{2n}\).  For a finitely generated fractional support, store exponents as
exact rationals and process them in increasing order.

\subsection{Specialized Gamma algorithm}

\begin{algorithm}[Gamma coefficient blocks]
\label{alg:gamma}
Let \(a_k\) be given by \cref{eq:gamma-ak}.  The following
procedure computes all blocks through \(d_{2N-1}\):
\begin{lstlisting}
E := 0
for n := 1,...,N:
    H := (1+E) log(1+E) - E
    for k := 1,...,n:
        H := H + a_k q^(2k) (1+E)^(-(2k-1))
    e_n := - coefficient(truncate_q(H,2n), q^(2n)) / Lambda
    E := E + e_n q^(2n)
return d_(2n-1) := e_n
\end{lstlisting}
All arithmetic is exact over \(\Q(\Lambda)\).  Bernoulli numbers are needed
only up to \(B_{2N}\).
\end{algorithm}

The recurrence is output-sensitive.  With naive dense series multiplication,
computing \(N\) blocks takes polynomial time in \(N\); FFT-based multiplication
can reduce the high-order cost.  Expression swell is controlled by retaining
factored powers of \(\Lambda\) and combining rational functions only at the
end of each step.

\subsection{Specialized Barnes algorithm}

\begin{algorithm}[Barnes coefficient blocks]
\label{alg:barnes}
Let \(b=-1/12\) and \(c_k\) be given by \cref{eq:barnes-ck}.  To compute
\(b_0,\ldots,b_N\):
\begin{lstlisting}
E := 0
for n := 1,...,N+1:
    H := K_lambda(E) - 2 lambda E
    H := H + 2 alpha q (1+E)
    H := H + 2 b q^2 (lambda+1+log(1+E))
    for k with 2k+2 <= n:
        H := H + 2 c_k q^(2k+2) (1+E)^(-2k)
    e_n := - coefficient(truncate_q(H,n),q^n) / (2 lambda)
    E := E + e_n q^n
return b_(n-1) := e_n
\end{lstlisting}
The parity check \(e_n(-\alpha)=(-1)^ne_n(\alpha)\) and the limiting check
\cref{eq:barnes-e-limits} are inexpensive regression tests.
\end{algorithm}

\subsection{Stable numerical evaluation}

For very large \(y\), never form \(\Gamma(x)\) or \(G(x)\) merely to evaluate
the inverse approximation.  Work with \(Y=\log y\) from the beginning.  The
recipes are:

\paragraph{Gamma.}
Compute
\[
 R=Y-\tfrac12\log(2\pi),\quad
 w=\W_0(R/\e),\quad T=R/w,\quad \Lambda=w+1,
\]
then evaluate \cref{eq:gamma-all-orders}.  The identity \(T=\e^{w+1}\) is
useful if \(R/w\) suffers loss of relative accuracy in a specialized numeric
regime, although the quotient is normally stable on the positive real axis.

\paragraph{Barnes.}
Compute
\[
 R=Y-\zeta'(-1),\quad
 v=\W_0(4R/\e^3),\quad
 U=2\sqrt{R/v},\quad \lambda=(v+1)/2,
\]
then evaluate \cref{eq:barnes-all-orders} and add one to recover
\(G_+^{-1}(y)\).

A single Newton correction in the logarithmic phase often supplies most or all
of the remaining working precision; a second correction retains quadratic
convergence when required.  For Gamma the stable update is
\begin{equation}
 x_{\mathrm{new}}=x-
 \frac{\log\Gamma(x)-Y}{\psi(x)},
 \label{eq:gamma-newton-log}
\end{equation}
where \(\psi\) is the digamma function.  This avoids overflow from evaluating
\(\Gamma(x)\) itself.  For Barnes \(G\), use the analogous update with the
logarithmic derivative of \(G\), or a high-accuracy numerical derivative of
\(\log G\) when no dedicated routine is available.

\subsection{Truncation policy}

At fixed moderate argument, adding terms indefinitely is counterproductive.
A robust policy is:
\begin{enumerate}
\item generate blocks in increasing \(q\)-order;
\item estimate the magnitude of the next block using the exact coefficient
      function, not only its leading large-\(v\) term;
\item stop near the smallest block, or sooner if the desired precision has
      been reached;
\item certify by evaluating the forward logarithmic residual and applying
      \cref{eq:mean-value-certificate} or one Newton step.
\end{enumerate}
Because the Lambert core is retained exactly, far fewer terms are usually
needed than in a fully expanded iterated-logarithmic formula.

\section{Formal, asymptotic, and convergent statements}
\label{sec:logical-status}

The calculations in this article use three notions of expansion.  Their
logical status should remain explicit.

\subsection{Formal Hahn identities}

Equations such as \cref{eq:gamma-normalized,eq:barnes-normalized} define
formal solutions coefficient by coefficient.  Their uniqueness is algebraic
and does not depend on an analytic function realizing the forward series.
Coefficient formulas, support restrictions, parity, and denominator bounds
live at this level.

\subsection{Poincar\'e asymptotic realization}

The forward Gamma and Barnes expansions are Poincar\'e asymptotic in sectors.
For every fixed truncation order, substitution of the formal inverse
truncation leaves a controlled residual.  Eventual monotonicity on the positive
real axis and \cref{thm:residual-certificate} then imply the inverse remainder.
This is the meaning of \cref{eq:gamma-all-orders,eq:barnes-all-orders}.

\subsection{Convergent local inverses}

The small positive Gamma branch in \cref{eq:gamma-small-lagrange} is an
ordinary local analytic inverse and therefore has a nonzero radius of
convergence in \(s=1/y\).  The large-branch Bernoulli expansions are not of
this kind.  The exact special-function inverse exists analytically on its
branch, but its displayed algebraic transseries is asymptotic rather than a
convergent sum of all Bernoulli blocks.

\subsection{Hyperasymptotic completion}

The late Bernoulli coefficients encode exponentially small remainders and
Stokes transitions.  Forward exponentially improved expansions for Gamma and
Barnes \(G\) have been developed by Nemes and earlier authors
\cite{NemesGamma,NemesBarnes}.  Applying
\cref{eq:inverse-exponential-sector} produces the leading inverse sectors.
A complete resurgent treatment would additionally track alien derivatives or
terminant multipliers through composition.  Nothing in the algebraic grid
prevents this extension, but its natural coefficient algebra is larger than
the one required for the explicit results here.

\section{Conclusions}
\label{sec:conclusions}

The large inverse branches of Gamma and Barnes \(G\) have a common normal
form.  Taking logarithms exposes a superlinear phase; a Lambert \(W\)
expression inverts its \(x^p\log x\) core exactly; a relative displacement
then satisfies an implicit equation with a unit linear coefficient.  The
remaining inversion is triangular in a Hahn grid.

For Gamma, the half-shift \(x=t+1/2\) removes every even inverse power from the
forward Stirling phase.  The inverse displacement consequently has only odd
powers \(T^{-1},T^{-3},\ldots\), with coefficient polynomials in
\(\Lambda^{-1}\).  The recurrence \cref{eq:gamma-recurrence} computes them to
arbitrary order and the forward remainder transfers with one additional
factor \(1/\Lambda\).

For Barnes \(G\), the linear term creates the first displacement
\(-\alpha/\lambda\), while the logarithmic term is resonant.  Its surviving
large-slope limit sums to the elementary algebraic series
\(\sqrt{1+q^2/6}-1\).  This both explains the coefficient pattern and provides
a strong symbolic check.  The actual inverse \(G_+^{-1}\) is obtained by
restoring the shift \(+1\).

The general calculus applies whenever the dominant phase can be normalized to
\(a x^p(\log x-\beta)\), or more generally to
\(a x^p(\log x-\beta)^r\), and all unabsorbed algebraic corrections have
positive relative power drop.  Its essential ingredients are exact core
inversion, a well-ordered support semigroup, a unit normalized slope, and a
residual certificate.  These conditions are constructive: they yield
coefficient algorithms rather than merely existence statements.

\begin{summarybox}
For practical use, keep Lambert \(W\) unexpanded, center the forward phase so
that its support is as sparse as possible, compute the inverse blocks by the
triangular recurrence, and validate the chosen truncation by its logarithmic
residual.  Expand into iterated logarithms only as a final presentation layer.
\end{summarybox}

\appendix

\section{Coefficient extraction identities}
\label{app:coefficient-extraction}

This appendix records two identities useful for independent implementations.

\subsection{Kernel coefficients}

From \cref{eq:kappa-general},
\begin{equation}
 \kappa_n(p,v)=\frac{v}{p}\binom pn+
 \binom pn\sum_{j=0}^{n-1}\frac{1}{p-j},
 \label{eq:kappa-harmonic-form}
\end{equation}
whenever none of \(p,p-1,\ldots,p-n+1\) is zero.  The right-hand side extends
polynomially through removable singularities because it is simply
\(\frac{v}{p}\binom pn+\partial_p\binom pn\).  For integral \(p\), it is safer
to evaluate the derivative form before substitution.

For Gamma, \(p=1\) and the universal nonlinear part simplifies to
\begin{equation}
 (1+E)\log(1+E)-E
 =\sum_{m\ge2}\frac{(-1)^m}{m(m-1)}E^m.
 \label{eq:gamma-kernel-series}
\end{equation}
For Barnes, direct expansion gives
\begin{equation}
 \Kcal_\lambda(E)
 =2\lambda E+\lambda E^2
 +E^2+\frac13E^3-\frac1{12}E^4
 +\frac1{30}E^5-\cdots.
 \label{eq:barnes-kernel-series}
\end{equation}
The separation of \(2\lambda E\) in the recurrence is essential; the term
\(\lambda E^2\) remains nonlinear and is responsible for the algebraic
resonant limit.

\subsection{Bell-polynomial form of the small Gamma branch}

Let
\begin{equation}
 \log\Gamma(1+x)=\sum_{m\ge1}\ell_m\frac{x^m}{m!},
 \qquad
 \ell_1=-\gamma,
 \qquad
 \ell_m=(-1)^m(m-1)!\zeta(m)\quad(m\ge2).
 \label{eq:ellm-loggamma}
\end{equation}
If \(Y_n\) denotes the complete exponential Bell polynomial, then
\begin{equation}
 \left.
 \frac{\dd^{n-1}}{\dd x^{n-1}}\Gamma(1+x)^n
 \right|_{x=0}
 =Y_{n-1}(n\ell_1,n\ell_2,\ldots,n\ell_{n-1}).
 \label{eq:small-gamma-bell}
\end{equation}
Consequently,
\begin{equation}
 \Gamma_0^{-1}(y)
 =\sum_{n\ge1}
 \frac{Y_{n-1}(n\ell_1,\ldots,n\ell_{n-1})}{n!\,y^n}.
 \label{eq:small-gamma-bell-series}
\end{equation}
This gives exact symbolic coefficients to arbitrary order without performing
series composition explicitly.

The next coefficient after \cref{eq:gamma-small-explicit} is
\begin{align}
 [y^{-6}]\,\Gamma_0^{-1}(y)=-\biggl({}&
 \frac{54\gamma^5}{5}+3\gamma^3\pi^2+6\gamma^2\zeta(3)
 +\frac{17\gamma\pi^4}{120}
 \notag\\
 &+\frac{\pi^2\zeta(3)}6+\frac{\zeta(5)}5
 \biggr).
 \label{eq:small-gamma-sixth}
\end{align}

\section{The sixth displayed Barnes block}
\label{app:barnes-b5}

For completeness, the coefficient multiplying \(U^{-5}\) in
\cref{eq:barnes-all-orders} is
\begin{equation}
 b_5(\lambda)=\frac{P_5(\lambda,\alpha)}{362880\lambda^{11}},
 \label{eq:barnes-b5}
\end{equation}
where
\begin{align}
 P_5={}&
 22680\alpha^6\lambda^5+40824\alpha^6\lambda^4
 -349272\alpha^6\lambda^3-352800\alpha^6\lambda^2
 \notag\\
 &+1111320\alpha^6\lambda-476280\alpha^6
 +11340\alpha^4\lambda^7+34020\alpha^4\lambda^6
 +3780\alpha^4\lambda^5
 \notag\\
 &-138600\alpha^4\lambda^4-132300\alpha^4\lambda^3
 +132300\alpha^4\lambda^2
 +1890\alpha^2\lambda^9
 \notag\\
 &+8568\alpha^2\lambda^8+15120\alpha^2\lambda^7
 +11718\alpha^2\lambda^6-3150\alpha^2\lambda^5
 -9450\alpha^2\lambda^4
 \notag\\
 &+105\lambda^{11}-668\lambda^{10}-378\lambda^9
 -21\lambda^8+175\lambda^7+105\lambda^6.
 \label{eq:barnes-P5}
\end{align}
The coefficient satisfies the predicted even parity in \(\alpha\), and
\begin{equation}
 \lim_{\lambda\to\infty}b_5(\lambda)
 =\frac{105}{362880}=\frac1{3456},
 \label{eq:barnes-b5-limit}
\end{equation}
matching \cref{eq:barnes-b-limits}.

\section{Direct residual checks}
\label{app:residual-checks}

The coefficient formulas can be checked without evaluating either inverse
function.  For Gamma, substitute
\begin{equation}
 E_N(q)=\sum_{n=1}^{N}d_{2n-1}(\Lambda)q^{2n}
 \label{eq:gamma-residual-EN}
\end{equation}
into the left-hand side of \cref{eq:gamma-normalized}.  The defining recurrence
is equivalent to
\begin{equation}
 \Lambda E_N+(1+E_N)\log(1+E_N)-E_N
 +\sum_{k=1}^{N}a_kq^{2k}(1+E_N)^{-(2k-1)}
 =O(q^{2N+2}).
 \label{eq:gamma-residual-formal}
\end{equation}
The coefficient of \(q^{2N+2}\) predicts the next block before any special
function is evaluated.

For Barnes, put
\begin{equation}
 E_N(q)=\sum_{n=1}^{N+1}b_{n-1}(\lambda)q^n.
 \label{eq:barnes-residual-EN}
\end{equation}
Then the left-hand side of \cref{eq:barnes-normalized}, with the Bernoulli sum
truncated wherever its exponent does not exceed \(N+1\), is
\begin{equation}
 O(q^{N+2}).
 \label{eq:barnes-residual-formal}
\end{equation}
These identities are exact in the rational-function coefficient field and
serve as stronger symbolic tests than decimal comparison.

For analytic verification at a numeric target, evaluate instead
\begin{equation}
 \delta_\Gamma=\log\Gamma(\widehat x)-\log y,
 \qquad
 \delta_G=\log G(\widehat x)-\log y,
 \label{eq:numeric-log-residuals}
\end{equation}
and divide by the corresponding logarithmic derivative.  This avoids the
severe scaling and overflow that would result from subtracting the original
function values.

\section{Notation dictionary}
\label{app:notation}

\begin{longtable}{>{\raggedright\arraybackslash}p{0.20\textwidth}
                  >{\raggedright\arraybackslash}p{0.72\textwidth}}
\toprule
Symbol & Meaning\\
\midrule
\endfirsthead
\toprule
Symbol & Meaning\\
\midrule
\endhead
\(Y\) & Logarithmic target, usually \(Y=\log y\).\\
\(C\) & Constant term removed from the forward phase.\\
\(R=Y-C\) & Target for the exact dominant core.\\
\(p,a,\beta\) & Parameters of \(a x^p(\log x-\beta)\).\\
\(v\) & Lambert variable \(\W((pR/a)\e^{-p\beta})\).\\
\(X\) & Exact Lambert core \(\e^\beta\e^{v/p}\).\\
\(q=X^{-1}\) & Exponentially small grid monomial relative to \(v^{-1}\).\\
\(\epsilon\) & Relative inverse displacement, \(x=X(1+\epsilon)\).\\
\(T,\Lambda\) & Gamma core and slope parameter: \(T=R/\W(R/\e)\),
                  \(\Lambda=\log T\).\\
\(U,\lambda\) & Barnes core and half-slope parameter:
                  \(U=2\sqrt{R/v}\), \(\lambda=(v+1)/2\).\\
\(d_{2n-1}\) & Gamma displacement coefficient multiplying \(T^{-(2n-1)}\).\\
\(b_n\) & Barnes displacement coefficient multiplying \(U^{-n}\).\\
\(C_G\) & Barnes constant \(\zeta'(-1)=1/12-\log A\).\\
\bottomrule
\end{longtable}

\begin{thebibliography}{99}

\bibitem{Barnes1900}
E.~W. Barnes,
\newblock The theory of the \(G\)-function,
\newblock \emph{Quarterly Journal of Pure and Applied Mathematics}
\textbf{31} (1900), 264--314.

\bibitem{BorweinCorless}
J.~M. Borwein and R.~M. Corless,
\newblock Gamma and factorial in the Monthly,
\newblock \emph{American Mathematical Monthly} \textbf{125} (2018),
400--424.

\bibitem{CorlessEtAl}
R.~M. Corless, G.~H. Gonnet, D.~E.~G. Hare, D.~J. Jeffrey, and D.~E. Knuth,
\newblock On the Lambert \(W\) function,
\newblock \emph{Advances in Computational Mathematics} \textbf{5} (1996),
329--359,
\newblock \url{https://doi.org/10.1007/BF02124750}.

\bibitem{DLMF511}
NIST Digital Library of Mathematical Functions,
\newblock \S5.11, Asymptotic expansions of the Gamma function,
\newblock \url{https://dlmf.nist.gov/5.11}, accessed 3 September 2026.

\bibitem{DLMF517}
NIST Digital Library of Mathematical Functions,
\newblock \S5.17, Barnes' \(G\)-function,
\newblock \url{https://dlmf.nist.gov/5.17}, accessed 3 September 2026.

\bibitem{FolitseJeffreyCorless}
K.~A. Folitse, D.~J. Jeffrey, and R.~M. Corless,
\newblock Properties and computation of the functional inverse of Gamma,
\newblock in \emph{2017 19th International Symposium on Symbolic and Numeric
Algorithms for Scientific Computing}, IEEE, 2017, 63--66,
\newblock \url{https://doi.org/10.1109/SYNASC.2017.00020}.

\bibitem{GerholdHubalekTomovski}
S. Gerhold, F. Hubalek, and \v{Z}. Tomovski,
\newblock Asymptotics of some generalized Mathieu series,
\newblock \emph{Mathematica Scandinavica} \textbf{126} (2020), 424--450,
\newblock \url{https://doi.org/10.7146/math.scand.a-121106}.

\bibitem{JeffreyWatt}
D.~J. Jeffrey and S.~M. Watt,
\newblock The inverse of the complex Gamma function,
\newblock arXiv:2311.16583, 2023,
\newblock \url{https://arxiv.org/abs/2311.16583}.

\bibitem{NemesGamma}
G. Nemes,
\newblock Error bounds and exponential improvements for the asymptotic
expansions of the Gamma function and its reciprocal,
\newblock \emph{Proceedings of the Royal Society of Edinburgh Section A:
Mathematics} \textbf{145} (2015), 571--596,
\newblock \url{https://doi.org/10.1017/S0308210513001558}.

\bibitem{NemesBarnes}
G. Nemes,
\newblock Error bounds and exponential improvement for the asymptotic
expansion of the Barnes \(G\)-function,
\newblock \emph{Proceedings of the Royal Society A} \textbf{470} (2014),
20140534,
\newblock \url{https://doi.org/10.1098/rspa.2014.0534}.

\bibitem{Olver}
F.~W.~J. Olver,
\newblock \emph{Asymptotics and Special Functions},
\newblock A K Peters, Wellesley, MA, 1997 reprint of the 1974 edition.

\bibitem{ParisKaminski}
R.~B. Paris and D. Kaminski,
\newblock \emph{Asymptotics and Mellin--Barnes Integrals},
\newblock Cambridge University Press, Cambridge, 2001.

\bibitem{Pedersen}
H.~L. Pedersen,
\newblock Inverses of Gamma functions,
\newblock \emph{Constructive Approximation} \textbf{41} (2015), 251--267.

\bibitem{Uchiyama}
M. Uchiyama,
\newblock The principal inverse of the Gamma function,
\newblock \emph{Proceedings of the American Mathematical Society}
\textbf{140} (2012), 1343--1348,
\newblock \url{https://doi.org/10.1090/S0002-9939-2011-11023-2}.

\bibitem{vanDerHoeven}
J. van der Hoeven,
\newblock \emph{Transseries and Real Differential Algebra},
\newblock Lecture Notes in Mathematics 1888, Springer, Berlin, 2006.

\bibitem{WhittakerWatson}
E.~T. Whittaker and G.~N. Watson,
\newblock \emph{A Course of Modern Analysis}, 4th ed.,
\newblock Cambridge University Press, Cambridge, 1927.

\end{thebibliography}

\end{document}
